\documentclass[12pt]{article}
\usepackage{fullpage}
\usepackage{amssymb, amsmath}
\usepackage{graphicx}
\usepackage{bbm}
\usepackage{appendix}

\usepackage[section]{placeins}
\makeatletter
\AtBeginDocument{%
\expandafter\renewcommand\expandafter\subsection\expandafter{%
\expandafter\@fb@secFB\subsection
}%
}
\makeatother

\newcommand{\argmax}{\operatorname*{arg \ max}}
\newcommand{\argmin}{\operatorname*{arg \ min}}
\newcommand{\minim}{\operatorname*{minimize}}
\usepackage{algorithm,algorithmic}
\usepackage{setspace}

\usepackage{blindtext}
\usepackage[table]{xcolor}
\usepackage{xr}
\usepackage{booktabs}
\usepackage{breqn}
\usepackage[
backend=biber,
giveninits=true,
uniquename=false,
date=year,
style=authoryear-comp,
uniquelist=false,
sorting=nty,
url=false,
isbn=false,
doi=false,
eprint=false,
maxbibnames=1,
minbibnames=1,
maxcitenames=1,
mincitenames=1
]{biblatex}
\usepackage{tikz}
\usepackage{subfigure}
\usepackage{array}
\usepackage{enumerate}
\usepackage{amsthm}
\usepackage[shortlabels]{enumitem}
\newcolumntype{C}[1]{>{\centering\arraybackslash}p{#1}}
\usetikzlibrary{shapes.geometric}
\usetikzlibrary{arrows,automata}
\usepackage{soul}
\newtheorem{alg}{Algorithm}
\newtheorem{prop}{Proposition}
\newtheorem{assumption}{Assumption}
\newtheorem{corol}{Corollary}
\newtheorem{defin}{Definition}
\newtheorem{theorem}{Theorem}
\newtheorem{lemma}{Lemma}
\DeclareMathOperator{\Y}{\mathbf{Y}}
\DeclareMathOperator{\G}{\mathbf{G}}
\DeclareMathOperator{\C}{\mathbf{C}}
\DeclareMathOperator{\E}{\mathbf{E}}
\DeclareMathOperator{\K}{\mathbf{K}}
\DeclareMathOperator{\Hn}{\mathbf{H}}
\DeclareMathOperator{\I}{\mathbf{I}}
\DeclareMathOperator{\bSigma}{\mathbf{\Sigma}}
\DeclareMathOperator{\vect}{vec}
\DeclareMathOperator{\tr}{tr}
\DeclareMathOperator{\trace}{tr}
\DeclareMathOperator{\diag}{diag}
\DeclareMathOperator{\Var}{Var}
\DeclareMathOperator{\cor}{Corr}
\DeclareMathOperator{\median}{median}
\DeclareMathOperator{\sd}{SD}
\DeclareMathOperator{\se}{SE}
\DeclareMathOperator{\cov}{cov}

\DeclareNameAlias{sortname}{last-first}

\addbibresource{mvREHE.bib}
\addbibresource{mvREHE_2.bib}



% Keywords command
\providecommand{\keywords}[1]
{
\small	
\textbf{\textit{Keywords---}} #1
}

\renewcommand\appendixpagename{Supplementary Materials}


% \title{Disentangling Genetic Contributions to Human Brain Connectivity Using an Efficient Estimator of Genetic Covariance}

\title{Genetic Regression Analysis  of Human Brain Connectivity Using an Efficient Estimator of Genetic Covariance}


\author{Keshav Motwani, Ali Shojaie, Ariel Rokem, and Eardi Lila}


\date{}{}

\newcommand{\as}[1]{\textcolor{cyan}{#1}}
\newcommand{\km}[1]{\textcolor{orange}{#1}}
\newcommand{\el}[1]{\textcolor{blue}{#1}}
\newcommand{\ariel}[1]{\textcolor{red}{Ariel: #1}}

\begin{document}
\def\spacingset#1{\renewcommand{\baselinestretch}%
{#1}\small\normalsize} \spacingset{1}

\maketitle

\vspace{-6mm}

\begin{abstract}


\noindent Non-invasive measurements of the human brain using magnetic resonance imaging (MRI) have significantly improved our understanding the brain's network organization by enabling measurement of anatomical connections between brain regions (structural connectivity) and their coactivation (functional connectivity).
Heritability analyses have established that genetics account for considerable intersubject variability in structural and functional connectivity. However, characterizing how genetics shape the relationship between structural and functional connectomes remains challenging, since this association is obscured by unique environmental exposures in observed data. To address this, we develop a regression analysis framework that enables characterization of the relationship between latent genetic contributions to structural and functional connectivity. Implementing the proposed framework requires estimating genetic covariance matrices in multivariate random effects models, which is computationally intractable for high-dimensional connectome data using existing methods. We introduce a constrained method-of-moments estimator that is several orders of magnitude faster than existing methods without sacrificing estimation accuracy. For the genetic regression analysis, we develop regularized estimation approaches, including ridge, lasso, and tensor regression. Applying our method to Human Connectome Project data, we find that functional connectivity is moderately predictable from structure at the genetic level (max $R^2
= 0.34$), though it is not directly predictable in the observed data (max $R^2 = 0.03$). This stark contrast suggests that unique environmental factors mask strong genetically-encoded structure-function relationships.
\end{abstract}




% while the functional connectome is not predictable from the structural connectome in the observed data, we find that it is predictable when restricted to the estimated genetic component, suggesting close relationships at the genetic level which are weakened by distinct environmental exposures. In addition, we characterize the leading directions of variation in the genetic and environmental components of the connectome. Critical to conducting the proposed analysis is the introduction of 
%, which constrains the solution to the space of positive semi-definite matrices. 

\begin{keywords}
\em Structural and functional connectivity, variance components, genetic covariance, genetic regression analysis
\end{keywords}

\spacingset{1.15} %\spacingset{1.9} % DON'T change the spacing!

\vspace{-3.5mm}

\section{Introduction}


Striking levels of intersubject heterogeneity have been observed in the structure and function of the human brain, even within narrowly defined groups of young and healthy subjects \parencite{benson2022variability}. Brain connectivity is no exception, suggesting that a large variety of connectivity configurations can result in cognitively healthy subjects \parencite{finn2015functional}. Two types of connectivity are commonly studied: \emph{structural connectivity}, which describes anatomical connections between brain regions through white matter fibers, and \emph{functional connectivity}, which measures the degree of co-activation of brain regions. The set of connections between pairs of macroscopic brain regions is referred to as the connectome.
% Two types of connectivity are commonly studied connectivity: \emph{structural connectivity}, which describes anatomical connections between brain regions through white matter fibers diffusion-weighted MRI (dMRI), which is sensitive to the trajectory of large bundles of axons that connect distant regions through the white matter, is used to assess the  between brain regions. Functional MRI (fMRI), which is sensitive to the blood oxygenation level dependent (BOLD) signal -- a correlate of neural activity, because it fluctuates with the activity-dependent metabolic demand in brain tissue -- is used to assess \emph{functional connectivity}. This is usually achieved by observing the degree of co-fluctuation of the BOLD signal in different brain regions. The set of connections between pairs of macroscopic brain regions is referred to as the connectome.

Structural connectivity is believed to provide the physical substrate supporting functional organization. This structure–function coupling has been extensively characterized \parencite{baum2020development, pan2025mapping, fotiadis2024structure}. However, while it is well established that both structural and functional connectivity are influenced by genetic \parencite{colclough2017heritability,elliott2018genomewide,gu2021heritability,bianco2023heritability} and environmental factors \parencite{kulick2020longterm,pini2023pollutome}, the specific aspects of structure–function coupling that are genetically rooted versus those shaped by environmental exposures are not yet well understood. Classical analyses of structure–function coupling in neuroimaging studies may mask underlying relationships when genetic and environmental factors contribute differently to this coupling. On the other hand, univariate heritability analyses are inadequate for addressing this question because they fail to capture the complex and multivariate nature of structure–function coupling.

In the present work, we developed a \textit{latent genetic regression analysis} framework that characterizes genetically encoded structure-function associations. Specifically, we represent observed connectomes as the sum of latent genetic and environmental components and model the genetic component of the functional connectome as a linear function of the genetic component of the structural connectome. To estimate this linear mapping, we develop a moment-matching genetic covariance estimator with a positive semidefinite constraint, which we use as a plug-in to estimate the linear regression coefficients within a ridge, lasso, or tensor regression framework. Our analysis reveals that functional connectomes show substantially higher predictability from structural connectomes at the genetic component level than at the phenotypic level, suggesting that strong genetically-encoded structure-function relationships exist but are masked by environmental heterogeneity. %Understanding these distinct genetic and environmental contributions to brain connectivity provides insights into the biological mechanisms underlying cognitive function.




% \km{Below are some sentences I've removed from the intro, with some comments from me. Please feel free to revert these changes!}


% The neuroimaging community has long been interested in linking these inter-individual variations in the connectome to behavioral, psychometric, and other factors, yet these models usually account for only a small part of the observed variability \parencite{smith2015positivenegative,xia2018linked,llera2019interindividual, mantwill2022brain, gell2023burden}. \km{Eardi, after reading more into this sentence to address Ali's comments, I'm not sure the relevance of this sentence here. It seems this is the opposite direction of what we want, how connectome explains variation in some other traits, rather than how we can explain variation in the connectome. Do we need to have this sentence in the intro? Please let me know if I am misinterpreting!} 


% The structural connectome is thought to provide
% the physical mechanisms that support the observed functional connections. Understanding more precisely how structural and functional connectivity relate is critical to gaining insight into how the brain's structural properties are linked to its physiological aspects \parencite{gu2021heritability}. \km{I've moved this to the end of Section 1.1 now. Is that okay?}



%We provide a more detailed overview of our data and proposed analysis in Section~\ref{sec:data}, but first, we review the literature on heritability analysis in the following section.

% Although genetic and environmental contributions to the connectomes at an individual level cannot be directly observed, our proposed analysis enables us to study which aspects of the relationship between structural and functional connectivity profiles are determined by genetic and environmental factors, as if these components were directly observable. 


\begin{figure}[!t]
\centering
\includegraphics[width=\textwidth]{overall_km.png}
\caption{Illustration of the proposed approach. Observed connectomes are modeled as the sum of latent genetic and environmental components. Our methodology characterizes the relationship between structural and functional connectivity within the genetic component through the proposed latent genetic regression analysis framework.}\label{fig:introfig}
\end{figure}

\vspace{-3.5mm}


\subsection{Motivating data and statistical analysis}\label{sec:data}

In this work, we use data from the Human Connectome Project (HCP), which contains multimodal imaging data from 1,200 young adults, including diffusion-weighted MRI for structural connectivity and resting-state functional MRI \parencite{vanessen2012human}. Critically, the dataset includes genetically verified family structures (149 monozygotic twin pairs, 94 dizygotic twin pairs), enabling identification of genetic and environmental contributions.


% TO DO describe kinship structure and ABCD model https://journals.plos.org/plosone/article/file?id=10.1371/journal.pone.0235860&type=printable


Figure~\ref{fig:introfig} illustrates our analysis. Structural connectomes are constructed from diffusion MRI using tractography to count streamlines connecting each pair of regions. Functional connectomes are constructed from temporal correlations of cortical activation between regions. As we describe in Section~\ref{subsec:datanalaysisresults}, when we perform classical linear regression to predict functional connectivity from structural connectivity in the observed data, we find weak associations (maximum observed $R^2$ = 0.03). This limited predictability motivates our investigation of whether stronger structure-function associations exist at the genetic level but are masked by environmental effects. We therefore model the genetic component of functional connectivity as a linear function of the genetic component of structural connectivity, with regression coefficients capturing intrinsic structure-function associations independent of environmental factors.


Implementing the proposed genetic regression analysis at connectome scale presents a major computational challenge. The analysis requires estimating the joint genetic covariance matrix of structural and functional connectomes, from which the regression parameters can be estimated. Under our modeling choices, the structural and functional connectomes together constitute a 4,624-dimensional phenotype, which existing likelihood-based methods cannot handle. We therefore developed a computationally efficient covariance estimator that renders the proposed analysis tractable and enables the use of high-dimensional regression frameworks for our genetic regression analysis.


% In Figure~\ref{fig:introfig}, we present an overview of our proposed analysis. 
% We use diffusion MRI and functional MRI to build subject-specific connectomes. In particular, a tractography reconstruction algorithm is first applied to diffusion MRI to generate streamlines, which are 3D curves that estimate the underlying white matter tracts. Using a cortical atlas that defines a set of regions of interest (ROI), these data are represented as a matrix that quantifies the strength of the structural connections between any pair of ROIs. In this work, we define the structural connectivity strength as the number of streamlines that connect any two regions. We also use functional MRI to compute the temporal correlation between the time series representing cortical activation within each ROI. This yields a measure of functional connectivity between any two ROIs, also represented with a matrix. As shown in the bottom panel of Figure~\ref{fig:introfig}, the main aim of this study is to disentangle and characterize the genetic and environmental contributions to the connectome. In particular, since the structural connectome is thought to provide the physical mechanisms that support the observed functional connections, we aim to understand more precisely how structural and functional connectivity relate, by disentangling the aspects of the structure-function coupling that arise independently through genetic and environmental mechanisms. This is critical to gaining insight into how the brain's structural properties are linked to its physiological functions \parencite{gu2021heritability}.

% The analysis outlined in Figure~\ref{fig:introfig} relies on estimating genetic and environmental covariances, which serve as sufficient statistics for the parameters of the proposed regression analysis. However, estimating these covariance matrices using existing methods is not computationally feasible for more than approximately 10 traits, whereas
% Our application involves the analysis of joint structural and functional connectomes using parcellations with 13 regions, amounting to 182 connections in total, and \km{we estimate it would take over 300 years if we used existing methods}. Therefore, analyzing such data requires new computationally efficient methods.
% our application involves analyzing joint structural and functional connectomes with 4624 connections (68 region parcellation of the brain). As discussed later, this scale is not tractable with existing methods. Therefore, analyzing such data requires new computationally efficient approaches that do not sacrifice estimation accuracy.


%For both of these aims, estimating the covariance matrix of the genetic and environmental components is required, meaning that existing univariate heritability models are not sufficient. However, existing approaches for estimation of variance components in multivariate linear mixed models are either extremely computationally burdensome or result in poor estimates with moderate sample sizes. Therefore, we propose a new computationally efficient method which matches the estimation accuracy of the existing computationally burdensome methods in many settings, as shown in our simulation studies. 


% We use a joint vectorized representation of structural and functional connectomes, and decompose these to study the covariance of connections in the genetic component of these connectomes. 

% Specifically, in the multivariate version of this model, the matrix of connectomes is assumed to be the sum of matrix-variate normally distributed components corresponding to genetics, common environment, and unique environment. These components are identified by their assumed row-covariance matrices. 

% First, we identify the main modes of variation in the genetic component of the connectome. Next, we identify which aspects of the genetic component of the structural and functional connectomes are most related to each other. Finally, we study the relationship between genetic contributions of anatomical connections in the brain to the genetic contributions of co-activation of brain regions. To ease visualization, the ROIs have been clustered into 13 communities, and the connectomes are accordingly depicted as $13 \times 13$ matrices. 

% The resulting classical PCA mode of variation reveals a coupled structural and functional increase or decrease in connectivity between any pair of regions, indicating a coherent direction of variation. However, the genetic principal mode of variation tells a more complex story. For instance, increased functional connectivity between the cingulo-opercular and auditory regions correlates with reduced connectivity within the auditory region itself, and vice-versa. Parallel findings are noted in the principal mode of variation for the environmental component. Further details regarding the results of our application can be found in Section \ref{sec:application}.

% The rest of the paper is organized as follows. In Section~\ref{sec:methods}, we present the proposed multivariate heritability model along with extensions to high-dimensional and functional data phenotypes. In Section 3, we empirically compare the runtime and estimation error of our method under different structures of the true covariance matrices to competing methods. In Section 4, we use the proposed method to characterize the genetic and environmental modes of variations in structural and functional connectivity, an analysis that would not be computationally feasible with existing methods. 
\vspace{-3.5mm}

\subsection{Related literature}\label{sec:litreview}

%I believe we should briefly mention here what the paper is about

% However, these models do not characterize the specific variations in these phenotypes that arise due to genetic and environmental factors.

%The parameters in the univariate model are scalar variances, and can be interpreted in terms of proportion of variance explained. Specifically, heritability of a trait is defined as proportion of variance of that trait explained by the genetic component in the model.

%The multivariate version of the mixed effects model additionally models the covariance between traits in the random effects corresponding to the genetic and environmental components. 

Our proposed approach bridges methods for heritability analysis and estimation of variance components. Therefore, in this section we review relevant literature from these fields.
Heritability analyses assess the extent to which trait variation is attributable to genetic factors, typically using additive genetics and common environment (ACE) models \parencite{maes2005ace}. These models are an application of linear mixed models, which partition the total phenotypic variance by including random effects for additive genetic (A), common environment (C), and unique environment (E) components, each defined by its expected covariance structure. Restricted maximum likelihood (REML) is the gold standard for variance component estimation, but it is computationally prohibitive for multivariate traits. Method-of-moments estimators are faster \parencite{lindquist2012estimating} but may produce non-positive estimates. Recently, constrained method-of-moments-based estimators that guarantee nonnegative variance estimates have been proposed and studied in \textcite{zou2017covariance} and \textcite{yue2021rehe}; these methods improve upon estimation performance of their unconstrained variants while matching their computational efficiency.  

% Most genetic studies of brain connectivity use massive univariate analyses, estimating heritability separately for each connection \parencite{colclough2017heritability,ge2017heritability,elliott2018genomewide}. While multidimensional extensions exist \parencite{tian2023bayesian, zhao2022bayesian, zhao2022genetic, risk2021ace, luo2019fsem}, they focus on heritability estimation rather than genetic covariances needed for regression analysis, or rely on suboptimal truncation approaches for positive semi-definite constraints.


% Heritability analyses are commonly used in genetics to assess the extent to which the observed variation in a univariate trait is attributable to genetic factors. 
% In order to assess the extent of observed variation in a univariate trait attributable to genetic factors, so called ``heritability analyses'' have become popular in the genetics community  \textcolor{red}{insert ref}. 
% These analyses typically utilize an additive genetics and common environment model \parencite{maes2005ace}, which assumes that the observed trait is the sum of independent univariate contributions from genetic, common environment, and unique environment components. The heritability of the trait is then defined as the proportion of variance explained by the genetic component. 

% To estimate heritability, the variance of each latent component model must be estimated. Restricted maximum likelihood (REML) is widely regarded the gold standard for variance component estimation. However, it can be computationally prohibitive, even for a single trait. Method-of-moments-based estimators are much faster \parencite{lindquist2012estimating}, but are not guaranteed to result in nonnegative estimates of variance . Recently, constrained method-of-moments-based estimators, which guarantee nonnegative variance estimates, have been proposed and studied in \textcite{zou2017covariance} and \textcite{yue2021rehe}; these methods improve upon estimation performance of their unconstrained variants while matching their computational efficiency.  

Previous work on structure-function coupling \parencite{gu2021heritability} investigated the role of genetics by first estimating subject-specific structure-function correlations and then assessing the heritability of this univariate trait. This approach enables the use of univariate heritability models, however, it does not directly model population-level relationships between genetic components of structural and functional connectivity.

Genetic contributions to multivariate traits, such as entire connectomes, have been typically investigated using massive univariate analyses, in which heritability is estimated separately for each trait. For example, numerous heritability studies of brain connectivity have analyzed one brain connection at a time using a univariate model; see, e.g., \textcite{colclough2017heritability,ge2017heritability} and \textcite{elliott2018genomewide}. %Similarly, the multivariate heritability measure proposed by \textcite{ge2016multidimensional} is a weighted average of the individual univariate heritabilities, and they utilize REML to estimate the heritability of each of the traits separately and average these estimates to obtain an estimate of multivariate heritability. 
Extensions to multidimensional imaging phenotypes have recently been proposed \parencite{tian2023bayesian, zhao2022bayesian, zhao2022genetic}. While these approaches improve estimation accuracy by borrowing information from neighboring imaging regions, they mostly focus on estimating heritability and not the genetic and environmental cross-covariance between distinct regions. Other multidimensional extensions are specifically designed for twin study designs \parencite{risk2021ace, luo2019fsem}. These methods either rely on a truncation approach to constrain the covariance estimate to be positive semi-definite, which may be suboptimal, or require estimating the covariance rank from the scree plot, a non-trivial task.

Massively univariate analyses do not estimate the covariance matrices of the genetic component, and thus do not provide any biological insight into the multidimensional structure of the genetic contributions. Since several multivariate statistical estimands are functions of the population covariance matrix,  estimating the entire \textit{covariance matrix} of the genetic component enables a large class of downstream analyses on the latent genetic component even though it cannot be directly observed. Although we focus here on linear regression to characterize the genetic structure-function coupling, our framework also facilitates other methods, such as principal component analysis and canonical correlation analysis.

%Massively univariate analyses estimate the variance of the genetic component for each of the traits. While this is useful for estimating the heritability of the traits, it does not provide any biological insight into the multidimensional structure of the genetic contributions.  In contrast, estimating the entire \textit{covariance matrix} of the multivariate genetic component can provide deeper insight into the disentangled genetic %aspects of the intricate interactions between traits. Moreover, due to the central role of covariance matrices in multivariate data analysis, the estimate of the covariance matrix opens the door to a large class of downstream analyses on the disentangled genetic component, even though it cannot be directly observed. Examples of such analyses include linear regression, principal components analysis, and canonical correlation analysis, which are among the most widely used methods in multivariate data analysis with parameters defined in terms of the covariance matrix of the traits. 

Unlike their univariate counterparts, the estimation of the covariance matrix of latent multi-dimensional components has received limited attention. REML estimation of the covariance matrices using \texttt{lme4} \parencite{bates2013fast} is computationally prohibitive with as few as 10 traits and 1,000 subjects. \textcite{zhou2014efficient} proposed more efficient algorithms for computing the REML estimator, though their computational approach does not extend to models with more than two latent components, such as the additive genetics and common environment model considered in this work. Commonly used software packages, such as GCTA \parencite{yang2011gcta}, only support up to two traits. To get around the computational bottleneck associated with ML-based estimators, method-of-moments-based estimators have been proposed \parencite{ge2016multidimensional, liegeois2019resting, anderson2021heritability}. However, these estimators have poor estimation performance compared to REML, particularly when using familial data in the additive genetics and common environment model. This is due to the high similarity between the kinship and shared environment indicator matrices, which is especially relevant in the HCP data analyzed in this work \parencite{ge2016multidimensional}. Furthermore, these methods may produce covariance estimates that are not positive semi-definite, complicating downstream analyses of the latent components. 
%through the estimated covariance matrices.

% JUST REMOVED 03/07/2025
% In addition to the limitations of univariate approaches, studies that utilize multivariate models are primarily based on an extended notion of scalar heritability to multivariate traits. For example, \textcite{feng2020partitioning} analyzes the overall heritability of the connectome using the notion of heritability introduced in \textcite{ge2016multidimensional}. Notably, this notion of heritability does not involve the off-diagonal terms of the covariance matrices, and could thus be equivalently formulated in terms of multiple univariate models. In fact, \textcite{ge2016multidimensional} adopt this approach, allowing them to utilize REML for heritability estimation given its superior performance over the method-of-moments-based estimator for familial data.

%Therefore, focusing on heritability alone does not allow for maximal interpretation of the data through the multivariate model.
%This is analogous to what classical principal component analysis could estimate had we actually observed the genetic and environmental components. 


% However, univariate heritability models only yield a scalar estimate of heritability defined as the proportion of variation in a particular phenotype that can be attributed to genetic factors. This means that their application to complex multidimensional phenotypes, such as connectomes, can only happen in the form of massive univariate studies. Therefore, a 

% On the other hand, bivariate models have been introduced in \parencite{almasy1997bivariate, almasy1998multipoint, blangero1997multipoint}. A few methods have been proposed to estimate variance components from twin design studies \parencite{luo2019fsem,risk2021ace}. Methods that do not require individual-level genotype data but instead use GWAS summary statistics, have also been proposed \parencite{bulik-sullivan2015ld}, but these tend to be less accurate \parencite{bulik-sullivan2015atlas}. 

% However, as mentioned previously, classical methods for its estimation, like (restricted) maximum likelihood (REML), are  Due to this computational bottleneck, \textcite{ge2016multidimensional} proposed a multivariate extension of  Haseman-Elston regression, a generalized method of moments estimator, and it was applied to estimating heritability of the connectome in \textcite{feng2020partitioning}. However, this estimator does not guarantee positive semi-definite estimates of the covariance matrices
\vspace{-3.5mm}

\subsection{Our contributions}

We introduce a latent genetic regression analysis framework that characterizes genetically-encoded structure-function associations by modeling relationships between genetic components of structural and functional connectivity. Implementing this analysis at the connectome scale requires efficient estimation of genetic and environmental covariance matrices. To this end, we introduce a constrained method-of-moments estimator that guarantees positive semidefinite estimates, closely matches REML’s accuracy, and provides substantial computational speedups. For $q$ connections and $n$ samples, our method requires $\mathcal{O}(n^2 q^2)$ operations upfront followed by $\mathcal{O}(\min(n, q)^3)$ operations per iteration, compared to at least $\mathcal{O}(n^3 q^3)$ operations per iteration for REML \parencite{bates2015fitting}.

To implement the proposed genetic regression analysis, we developed regularized estimation approaches based on plug-in estimates of the genetic covariance matrix, including ridge regression, lasso regression, and tensor regression that leverages connectome structure. These methods estimate linear mappings from the latent genetic component of the structural connectome to the latent genetic component of the functional connectome, without requiring direct observation of the genetic component itself.

Applied to HCP data, we find functional connectivity is substantially more predictable from structural connectivity at the genetic level than in observed data (maximum $R^2$: 0.34 vs 0.02). This stark contrast suggests unique environmental factors mask genetically-encoded relationships, and the uncovered genetically-encoded relationships may provide insights into the biological architecture of structure-function coupling.

The remainder of the paper is organized as follows. Section~\ref{sec:methods} formulates the genetic regression framework and introduces our proposed genetic covariance estimator. Section~\ref{sec:application} applies our methodology to HCP data. Section~\ref{sec:simulations} presents extensive simulations validating our approach. We end with some concluding remarks in Section~\ref{sec:discussion}. Supplementary Materials contain the derivation of the numerical algorithm, asymptotic theory, additional simulations, and extensions to function-valued phenotypes.


% To perform the proposed regression analysis of the latent components, we introduce a generalized, constrained method-of-moments estimator that guarantees positive semi-definite covariance matrix estimates. The proposed method results in improved estimation performance, nearly matching that of REML, which is widely considered the gold standard for variance component estimation. We also propose a block-coordinate descent algorithm that can efficiently incorporate this positive semi-definite constraint and is many orders of magnitude faster than REML.  
% For $q$ connections and $n$ samples, our method requires one $\mathcal{O}(n^2 q^2)$ operation upfront by $\mathcal{O}(\min(n, q)^3)$ operations per iteration.
% .The runtime of our method is $\mathcal{O}(n^2q^2 + T q^3)$ 
% This is a significant improvement over the runtime of REML, for which an \texttt{lme4}-based implementation would involve operations scaling with at least $n^3 q^3$ in each iteration to simply evaluate the REML criterion \parencite{bates2015fitting}. The two-component model in \textcite{zhou2014efficient} has a runtime complexity of $\mathcal{O}(n^3 + T_1 n q^2 + T_2 n q^6)$, where $T_1$ and $T_2$ represent the number of expectation-maximization and Newton-Raphson iterations, respectively. Although directly comparing runtimes is not straightforward due to their dependency on these optimization parameters, simulations demonstrate that our proposed approach outperforms \textcite{zhou2014efficient} in speed without compromising estimation accuracy. This improved computational efficiency will enable the analysis of neuroimaging and other multivariate datasets as the parcellations get more refined and datasets grow larger, providing insights into the genetic and environmental mechanisms underlying multidimensional neurological traits, such as the connectome.


% The remainder of the paper is organized as follows. In Section~\ref{sec:methods}, we 
%review univariate heritability models and 
% formulate the proposed genetic and environmental regression analysis and introduce the constrained method-of-moments estimator for the genetic and environmental covariance matrices. In Section~\ref{sec:application}, we apply the proposed estimator to data from the HCP to characterize the relationship between functional and structural connectomes via genetic, common environmental, and unique environmental mechanisms and identify the dominant modes of variation in the genetic component. In Section~\ref{sec:simulations}, we conduct extensive simulations to validate the proposed latent covariance estimation method and latent regression analysis in settings that closely resemble our application. We end with some concluding remarks in Section~\ref{sec:discussion}. The derivation of the numerical algorithm, the proof of asymptotic convergence of the estimator, and additional simulation results are provided in the Supplementary Materials. To further demonstrate the flexibility of the proposed method, in the Supplementary Materials, we also introduce an extension for function-valued phenotypes, which takes advantage of their smoothness. This highlights one of several ways the proposed framework can be combined with other covariance estimation methods to incorporate additional prior information about the data.


%Give (again) an example of these numbers in out in our setting with (estimated) runtimes 




% In this work, we characterize the modes of variations in structural and functional connectivity of the human brain, that can be attributed to genetic or environmental factors, through a novel statistical model that allows for efficient estimation of the variance components of multidimensional phenotypes. This is orders of magnitude faster than existing methods and can be applied to simultaneously analyze hundreds of imaging phenotypes. Its multivariate nature allows the proposed model to provide insights into the intricate interactions network between structural and functional connectivity of different regions, disentangling multidimensional genetic and environmental contributions.

\vspace{-3.5mm}

\section{Methodology}\label{sec:methods}
\vspace{-3.5mm}



\subsection{Latent genetic regression analysis}

% Understanding the complex interplay of genetic and environmental factors in shaping brain connectivity is central to understanding brain function \parencite{gu2021heritability}. The complex network of connections, both structural and functional, is shaped by a combination of genetic predispositions and environmental exposures throughout development. Traditional statistical analyses examine these connections in their observed form, conflating genetic and environmental contributions.

%By disentangling these components in a statistical model, we can perform analyses on the ``pure'' genetic architecture of the brain -- examining relationships that might otherwise be masked by environmental influences. This approach allows us to investigate the underlying biological blueprint of neural connectivity. 
To disentangle genetic and environmental factors in shaping brain connectivity, we adopt a multivariate additive genetics and common environment model, in which a  $q$-dimensional trait $\mathbf{y}_i$ is modeled as the sum of independent genetic, common environment, and unique environment random components:
\begin{equation}  \label{eq:multivariateindivmodel2}
\mathbf{y}_i = \mathbf{g}_i + \mathbf{c}_i + \mathbf{e}_i
\end{equation}
where
$\mathbf{g}_i \sim \mathcal{N}(0, \bSigma_G)$, $\mathbf{c}_i \sim \mathcal{N}(0, \bSigma_C)$, and $\mathbf{e}_i \sim \mathcal{N}(0, \bSigma_E)$. Here, $\mathcal{N}(0, \bSigma)$ denotes the multivariate normal distribution with mean $0$ and covariance matrix $\bSigma$ and $\bSigma_G, \bSigma_C$ and $\bSigma_E$ are unknown $q \times q$ covariance matrices. Specifically, we model the vectorized structural and functional connectomes as a multivariate trait $\mathbf{y}_i = \left(\mathbf{y}_i^{\mathcal{S}}, \mathbf{y}_i^{\mathcal{F}} \right)$ and partition the latent components conformably, e.g., $\mathbf{g}_i = \left(\mathbf{g}_i^{\mathcal{S}}, \mathbf{g}_i^{\mathcal{F}} \right)$.

To characterize the association between structural and functional brain connections driven by genetic factors, our primary goal is to perform linear regression of the functional connectome on the structural connectome, restricted to the genetic component of each. Specifically, we fit separate regressions for each functional connection, using the full structural connectome as the predictor. Formally, for the $j$th functional connection, we are interested in the best linear approximation of its genetic component based on the genetic component of the entire structural connectome:
\begin{equation} \label{eq:geneticbeta}
\boldsymbol{\beta}_{G, j}^* = \argmin_{\boldsymbol{\beta}} \mathbb{E}\left[ \left\|\mathbf{g}_{i}^{\mathcal{F}_j} - \boldsymbol{\beta}^\top \mathbf{g}_{i}^{\mathcal{S}} \right\|_2^2\right]
\end{equation}
and the associated proportion of genetic variance explained (genetic-$R^2$):
\begin{equation} \label{eq:geneticr2}
R^2_{G, j}(\boldsymbol{\beta}) = 1 - \frac{\Var\left(\mathbf{g}_{i}^{\mathcal{F}_j} - \boldsymbol{\beta}^\top \mathbf{g}_{i}^{\mathcal{S}}\right)}{\Var\left(\mathbf{g}_{i}^{\mathcal{F}_j}\right)}.
\end{equation}

These genetic regression coefficients capture the intrinsic associations between structural and functional connectomes that exist independent of environmental influences, representing the genetic ``blueprint'' of structure-function relationships.

In previous studies of structure–function coupling \parencite{gu2021heritability}, the coupling is typically characterized by a subject-specific parameter—such as the correlation between the observed structural and functional connectomes—and the heritability of these parameters is reported. These approaches implicitly assume that each functional connection influences only its corresponding structural connection to enable the estimation of subject-specific correlations. In contrast, genetic $R^2$ captures more complex population-level relationships between structure and function. Additional details on the interpretation of these parameters in the context of structure–function coupling studies are left to Section~\ref{sec:application}.

Although the genetic regression and genetic $R^2$ parameters are defined in terms of latent unobservable variables, our proposed analysis leverages the observation that they can be expressed in terms of specific blocks of the genetic covariance matrix, $\bSigma_G$. We partition the genetic covariance matrix according to the structural and functional components:
$$
\bSigma_G = 
\begin{bmatrix}
\bSigma_G^{\mathcal{S},\mathcal{S}} & \bSigma_G^{\mathcal{S},\mathcal{F}} \\
\bSigma_G^{\mathcal{F},\mathcal{S}} & \bSigma_G^{\mathcal{F},\mathcal{F}}
\end{bmatrix}. 
$$
The genetic regression parameter can then be identified as
\begin{equation}\label{eq:geneticbeta_cov}
\boldsymbol{\beta}_{G, j}^*  = 
\argmin_{\boldsymbol{\beta}} \boldsymbol{\beta}^\top \bSigma_G^{\mathcal{S}, \mathcal{S}} \boldsymbol{\beta} - 2\boldsymbol{\beta}^\top \bSigma_G^{\mathcal{S},\mathcal{F}_j},
\end{equation}
and the genetic $R^2$ as
\begin{equation}\label{eq:geneticr2_cov}
R^2_{G, j}(\boldsymbol{\beta}) = 1 - \frac{\bSigma_G^{\mathcal{F}_j,\mathcal{F}_j} - 2 \boldsymbol \beta^\top \bSigma_G^{\mathcal{S}, \mathcal{F}_j} + \boldsymbol \beta^\top \bSigma_G^{\mathcal{S}, \mathcal{S}} \boldsymbol \beta }{\bSigma_G^{\mathcal{F}_j, \mathcal{F}_j}}.
\end{equation}

Therefore, estimating the genetic covariance matrix, $\bSigma_G$, enables plug-in estimation of the genetic regression parameters. The plug-in loss minimization framework naturally accommodates regularized approaches, including ridge and lasso regression, as well as extensions to tensor regression that respect the matrix structure of the structural connectome predictor. We first introduce our efficient covariance estimation methodology in the next section, then develop the latent regression estimation methodology in Section~\ref{sec:regression_estimation}.

\subsection{Estimation of covariance matrices}
\subsubsection{Observed data model}
%To estimate the covariance matrices $\bSigma_G$, $\bSigma_C$, and $\bSigma_E$ from observed data, we need to consider the model across multiple subjects simultaneously. 
Let $\mathbf{y}_1, \dots, \mathbf{y}_n$ be $n$ identically distributed observations according to the model in \eqref{eq:multivariateindivmodel2} and define $\Y \in \mathbb{R}^{n \times q}$ as the matrix with rows $\mathbf{y}_i^\top$. Similarly defining matrices $\G$, $\C$, and $\E$ for the genetic, common environmental, and unique environmental components across all subjects, we can write
\begin{equation} \label{eq:multivariatevariateobservedmodel}
\Y = \G + \C + \E.
\end{equation}
For an $n \times q$ matrix-valued random variable $\mathbf{Z}$, we say $\mathbf{Z} \sim \mathcal{MN}(\mathbf{0}, \mathbf{D}, \bSigma)$ follows a matrix-normal distribution with row-covariance $\mathbf{D}$ (describing dependence between subjects) and column-covariance $\bSigma$ (describing dependence between traits). This is formally defined as $\vect(\mathbf{Z}) \sim \mathcal{N}(0, \bSigma \otimes \mathbf{D})$, where $\vect(\cdot)$ denotes vectorization and $\otimes$ is the Kronecker product.

If observations were independent across subjects, we would have $\G \sim \mathcal{MN}(\mathbf{0}, \mathbf{I}_n, \bSigma_G)$, $\C \sim \mathcal{MN}(\mathbf{0}, \mathbf{I}_n, \bSigma_C)$, and $\E \sim \mathcal{MN}(\mathbf{0}, \mathbf{I}_n, \bSigma_E)$. However, since the row-covariances are the same, only the sum $\bSigma_G + \bSigma_C + \bSigma_E$ is identifiable from the data---the individual covariance matrices cannot be separately estimated.

Identifying the individual covariance matrices therefore requires an experimental design that includes related subjects with varying degrees of genetic relatedness and shared environmental exposure. Let $\K_n$ and $\Hn_n$ be $n \times n$ positive semi-definite matrices with diagonal entries equal to 1, where off-diagonal entries quantify the genetic relatedness and shared environment status between each pair of subjects, respectively:
\begin{equation}
\label{eq:multivariatevariateobservedmodel2}
\G \sim \mathcal{MN}(\mathbf{0}_{n \times q}, \K_n, \bSigma_G), \quad \C \sim \mathcal{MN}(\mathbf{0}_{n \times q}, \Hn_n, \bSigma_C), \quad \E \sim \mathcal{MN}(\mathbf{0}_{n \times q}, \I_n, \bSigma_E).
\end{equation}

Under the model in \eqref{eq:multivariatevariateobservedmodel} and \eqref{eq:multivariatevariateobservedmodel2}, the individual covariance matrices $\bSigma_G$, $\bSigma_C$, and $\bSigma_E$ become identifiable provided that $\K_n$, $\Hn_n$, and $\mathbf{I}_n$ are linearly independent, which is satisfied when the study includes subjects with varying degrees of genetic relatedness and shared environmental exposure.

Existing methods for estimation of the covariance matrices in the multivariate random effects model are computationally expensive, especially with a large number of traits. To overcome this challenge, we propose an extension of \textcite{yue2021rehe} to the multivariate setting, using a constrained method-of-moments estimator rather than maximum likelihood. Our approach requires only the first two moments of the data, specifically the Kronecker covariance structure: e.g. $\cov(\vect(\G))=\bSigma_G \otimes \K_n$. Therefore, while normal distributions are used for expositional clarity, the normality assumption is not required for our estimation procedure.

% $$
% Y = \sum_{k = 1}^K \Gamma_k + \epsilon
% $$
% $$
% Y \in \mathbb{R}^{n \times q}, \Gamma_k \sim \mathcal{MN}(0_{n \times q}, D_k, \bSigma_k), \epsilon \sim \mathcal{MN}(0_{n \times q}, I_n, \bSigma_e).
% $$
% Our goal is to estimate the variance components $\bSigma_k$ for $k = 1, \dots, K$ and $\bSigma_e$. 

% Without loss of generality, consider a simplified model with with only one random effect:
%  In our primary application of interest, $K_n$ is a matrix encoding known information about relatedness among subjects, so we refer to $G$ as the genetic component, and $E$ as the environmental component, explaining variation that is not shared between subjects.



%  and canonical correlation analysis on the genetic component between two sets of traits $\mathcal{X}$ and $\mathcal{Y}$ estimates parameters of the form
% \begin{equation}\label{eq:mode}
% \hat{\mathbf{w}_1}, \hat{\mathbf{w}_2} = \argmax_{\mathbf{w}_1, \mathbf{w}_2} \frac{\mathbf{w}_1^\top  [\bSigma_G]_{\mathcal{X}, \mathcal{X}}^{-1} [\bSigma_G]_{\mathcal{X}, \mathcal{Y}} [\bSigma_G]_{\mathcal{Y}, \mathcal{Y}}^{-1} \mathbf{w}_2}{\sqrt{\mathbf{w}_1^\top \mathbf{w}_1}\sqrt{\mathbf{w}_2^\top \mathbf{w}_2}}
% \end{equation} 


%In this setting, heritability is often defined as $h^2 = \tr(\bSigma_G) / [\tr(\bSigma_E) + \tr(\bSigma_G)]$. Note, however, that even though this model involves the covariances between columns of $\mathbf{G}$ and $\mathbf{E}$, and therefore the interaction between multiple traits, heritability only depends on their diagonal entries -- the variances. 
%Therefore, simply analyzing heritability does not fully take advantage of the multivariate model, and could be equivalently estimated by fitting a univariate model for each trait separately. However, analyzing $\bSigma_G$, the ``genetic covariance'' between traits, has been of particular interest recently \parencite{van2019genetic, ni2018estimation} to understand the shared action of genes on pairs of traits. 

\vspace{-3.5mm}

\subsubsection{Restricted multivariate Haseman-Elston regression}

The model introduced is a special case of the more general model with $K + 1$ components:
\begin{equation} \label{eq:genmodel}
\Y = \sum_{k = 0}^K \boldsymbol{\Gamma}_k,
\end{equation}
where $\boldsymbol{\Gamma}_k \sim \mathcal{MN}(\mathbf{0}_{n \times q}, \mathbf{D}_k, \bSigma_k)$. The model in (\ref{eq:multivariatevariateobservedmodel}) can be recovered by setting $\mathbf{D}_0 = \I_n$, $\mathbf{D}_1 = \K_n$ and $\mathbf{D}_2 = \Hn_n$ so that $\G = \boldsymbol{\Gamma_1}$, $\C = \boldsymbol{\Gamma_2}$, and $\E = \boldsymbol{\Gamma_0}$.
%and $\bSigma_0 = \bSigma_E, \bSigma_1 = \bSigma_G$, and $\bSigma_2 = \bSigma_C$ 

Classical Haseman-Elston (HE) regression \parencite{haseman1972investigation} estimates the variance components in a univariate model using a generalized method-of-moments estimator that minimizes the squared difference between theoretical and empirical second moments. To generalize this to the multivariate setting, we use the fact that the second moments of $\Y$ under the model in (\ref{eq:genmodel}) are given by
\begin{equation} \label{eq:ee}
\mathbb{E}[\Y_{i, j} \Y_{l, m}] = \sum_{k = 0}^K [\bSigma_k]_{j, m} [\mathbf{D}_k]_{i, l},
\end{equation}
for all $i,l \in [n], j,m \in [q]$. 
Therefore, a multivariate extension of HE regression would minimize
\begin{align} 
\hat \bSigma_0, \dots, \hat \bSigma_K &= \argmin_{\bSigma_k} \sum_{j = 1}^q \sum_{m = 1}^{q} \sum_{i = 1}^n \sum_{l = 1}^{n} \left(\Y_{i,j}\Y_{l,m} - \sum_{k = 0}^K [\bSigma_k]_{j, m} [\mathbf{D}_k]_{i, l} \right)^2 \label{eq:mvHE} \\
&= \argmin_{\bSigma_k} \sum_{j = 1}^q \sum_{m = 1}^{q} \left\|\Y_{:,j} \Y_{:,m}^\top   -  \sum_{k = 0}^K [\bSigma_k]_{j, m} \mathbf{D}_k \right\|_F^2. \label{eq:leastsquares}
\end{align}
The optimization problem \eqref{eq:leastsquares} can be solved for each pair of outcome variables $(j, m) \in [q] \times [q]$ separately, which can be rewritten as a least-squares linear regression. Specifically, let
$
\tilde{\Y}^{(j, m)} = \vect(\Y_{:, j} \Y_{:, m}^\top),
$
$
\tilde{\mathbf{X}} = (\vect(\mathbf{D}_0), \dots, \vect(\mathbf{D}_K)),
$
and
$
{\sigma}^{(j, m)} = ([\bSigma_0]_{j, m}, \dots, [\bSigma_K]_{j, m})^\top
$. Then solving (\ref{eq:leastsquares}) is equivalent to solving
\begin{equation*} \label{eq:min}
\hat \sigma^{(j, m)} = \argmin_{\sigma^{(j, m)}} \left\| \tilde{\Y}^{(j, m)} - \tilde{\mathbf{X}} {\sigma}^{(j, m)}\right\|_2^2
\end{equation*} with $[\hat{\bSigma}_k]_{j, m} = [\hat{\sigma}^{(j, m)}]_{k+1}$. This problem has a closed-form solution given by
$
\hat{\sigma}^{(j, m)} = (\tilde{\mathbf{X}}^{\top} \tilde{\mathbf{X}})^{-1} \tilde{\mathbf{X}}^{\top} \tilde{\Y}^{(j, m)}.
$


The estimator in \eqref{eq:leastsquares} has been proposed before by \textcite{ge2016multidimensional} for heritability analyses of neuroimaging traits and has been used by \textcite{liegeois2019resting} and \textcite{anderson2021heritability}. However, this estimator does not return positive semi-definite (PSD) estimates in general, requiring ad-hoc corrections before performing analyses using the covariance matrix. A simple correction, analogous to truncating the variance estimates at $0$ in the univariate case, is to truncate the eigenvalues of $\hat{\bSigma}_k$ at $0$. We refer to the truncated version of \eqref{eq:leastsquares} as the multivariate Haseman-Elston (mvHE) estimator.

Although truncating the estimate obtained from \eqref{eq:leastsquares} results in positive semi-definite estimates of the covariance matrices, these estimates do not minimize \eqref{eq:leastsquares} subject to the constraint that each of the $\bSigma_k$ is positive semi-definite (PSD). Thus, we consider solving \eqref{eq:leastsquares} with positive semi-definiteness constraints, which we call the restricted multivariate Haseman-Elston (mvREHE) estimator:
\begin{align}
\hat \bSigma_0, \dots, \hat \bSigma_K &= \argmin_{\bSigma_k \succeq 0} \sum_{j = 1}^q \sum_{m = 1}^{q} \left\|\Y_{:,j} \Y_{:,m}^\top   -  \sum_{k = 0}^K [\bSigma_k]_{j, m} \mathbf{D}_k \right\|_F^2 \label{eq:mvREHEunpen1} \\
&= \argmin_{\bSigma_k \succeq 0}  \sum_{i = 1}^n \sum_{l = 1}^{n} \left\|\Y_{i,:}^\top \Y_{l,:}   - \sum_{k = 0}^K [\mathbf{D}_k]_{i, l} \bSigma_k \right\|_F^2. \label{eq:mvREHEunpen}
\end{align}

The optimization problem in \eqref{eq:mvREHEunpen1} is no longer separable across pairs of outcome variables due to the PSD constraint. Therefore, to solve this optimization problem, we use a block coordinate descent algorithm. Specifically, we update $\bSigma_z$ with $\bSigma_k$ for $k \neq z$ held fixed at $\hat \bSigma_k^{(t)}$:
\begin{equation} \label{eq:update}
\hat \bSigma_z^{(t+1)} = \argmin_{\bSigma_z \succeq 0} \sum_{i = 1}^n \sum_{l = 1}^{n} \left\|\left(\Y_{i,:}^\top \Y_{l,:} - \sum_{k\neq z} [\mathbf{D}_k]_{i, l} \hat \bSigma_k^{(t)} \right) - [\mathbf D_k]_{i, l} \bSigma_z \right\|_F^2.
\end{equation}
Since the objective function in (\ref{eq:mvREHEunpen}) is convex and differentiable, with each of the blocks constrained to a convex set, block coordinate descent updates will converge to the solution of (\ref{eq:mvREHEunpen}).
The following lemma, proved in Appendix \ref{prooflemma}, allows us to compute each of these updates in closed form.
\begin{lemma} \label{lemma:update}
The solution to 
$$
\hat \bSigma = \argmin_{\bSigma \succeq 0}  \sum_{w = 1}^W \left\|\mathbf{S}_w - a_w \bSigma \right\|_F^2
$$
is $\hat \bSigma = \mathbf U \mathbf \Lambda^+ \mathbf U^\top$, where $\mathbf U \mathbf \Lambda \mathbf U^\top$ is the eigendecomposition of 
$$
\mathbf{S} = \left( \sum_{w = 1}^W a_w^2 \right)^{-1} \sum_{w = 1}^W a_w \mathbf{S}_w
$$
and $\mathbf \Lambda^+_{j,j} = \max(\mathbf \Lambda_{j,j}, 0)$.
\end{lemma}

Following Lemma \ref{lemma:update}, to solve (\ref{eq:update}), we must compute the eigendecomposition of
\begin{align*}
&\left( \sum_{i = 1}^n \sum_{l = 1}^n [\mathbf{D}_z]_{i, l}^2 \right)^{-1} \sum_{i = 1}^n \sum_{l = 1}^n [\mathbf{D}_z]_{i, l} \left(\Y_{i,:}^\top \Y_{l,:} - \sum_{k \neq z}[\mathbf{D}_k]_{i, l} \hat \bSigma_k^{(t)} \right) \\
&=\left( \underbrace{\sum_{i = 1}^n \sum_{l = 1}^n [\mathbf{D}_z]_{i, l}^2}_{\mathbf{Q}_{z, z}}  \right)^{-1} \left(\underbrace{\Y^\top \mathbf{D}_z \Y}_{\mathbf{W}_z} - \sum_{k \neq z} \hat \bSigma_k^{(t)} \underbrace{\sum_{i = 1}^n \sum_{l = 1}^n [\mathbf{D}_z]_{i, l} [\mathbf{D}_k]_{i, l}}_{\mathbf{Q}_{z, k}} \right).
\end{align*}
% \begin{align*}
% &\left( \sum_{i = 1}^n \sum_{l = 1}^n [\K_n]_{i, l}^2\right)^{-1} \sum_{i = 1}^n \sum_{l = 1}^n [\K_n]_{i, l} \left(\Y_{i,:}^\top \Y_{l,:} - [\I_n]_{i, l} \bSigma_E^{(t)} \right) \\
% &=\left( \underbrace{ \sum_{i = 1}^n \sum_{l = 1}^n [\K_n]_{i, l}^2}_{Q_{G}}\right)^{-1} \left(\underbrace{\sum_{i = 1}^n \sum_{l = 1}^n [\K_n]_{i, l} \Y_{i,:}^\top \Y_{l,:}}_{W_G} -  \bSigma_E^{(t)} \underbrace{\sum_{i = 1}^n \sum_{l = 1}^n [\K_n]_{i, l} [\I_n]_{i, l}}_{Q_{G, E}} \right)
% \end{align*}
To expedite the computation, we can pre-compute $\mathbf{Q}_{k,k'}$ and $\mathbf{W}_{k}$ for all $k, k' = 1, \dots K$ as they do not depend on $\hat \bSigma_k^{(t)}$. If $q \gg n$, the eigendecomposition may be computationally prohibitive. However, the following proposition allows us to compute eigendecompositions of $n \times n$ matrices rather than $q \times q$ matrices in each iteration, simply by computing the compact SVD, $\Y = \mathbf{UDV}^\top$, up front, and computing the mvREHE estimator with data matrix $\Y\mathbf{V} \in \mathbb{R}^{n \times n}$. Thus, in practice, we must only compute eigendecompositions of $\min(n, q) \times \min(n, q)$ matrices.

\begin{prop}
Let $\Y = \mathbf{UDV}^\top$ be the compact SVD of $Y$. Then, the solution to \eqref{eq:mvREHEunpen} is $\mathbf{V}\tilde \bSigma_k \mathbf{V}^\top$ where $\tilde \bSigma_k$ is the solution to \eqref{eq:mvREHEunpen} using $\mathbf{YV}$ in place of $\Y$. 
\label{prop:runtime}
\end{prop}

Algorithm~\ref{alg:original} summarizes the procedure. The runtime of our algorithm is only $\mathcal{O}(nq^2 + n^2 \min(q, n)^2 + TK\min(q, n)^3)$, where $T$ is the number of iterations. In comparison, an \texttt{lme4}-based implementation of REML would involve operations scaling with at least $n^3 q^3$ in each iteration to simply evaluate the REML criterion. This is because REML involves computing the Cholesky decomposition of an $nq \times nq$ matrix \parencite{bates2015fitting}. Even the optimized algorithm for REML proposed by \textcite{zhou2014efficient}, which only applies to models with 2 components, or $K = 1$ in \eqref{eq:genmodel}, has a runtime complexity of $\mathcal{O}(n^3 + T_1 n q^2 + T_2 n q^6)$, where $T_1$ is the number of expectation-maximization iterations and $T_2$ is the number of Newton-Raphson iterations. The computational speed-up of our method makes it very appealing for high-dimensional and large sample applications. 

\begin{algorithm}[b]\caption{Blockwise coordinate descent algorithm for solving (\ref{eq:mvREHEunpen})} \label{alg:original}
\setstretch{1.15}
\textcolor{white}{X} 1. For $z \in [K]$, compute $\mathbf W_z = \Y^\top \mathbf{D}_z \Y$  \\
\textcolor{white}{X} 2. Compute $\mathbf Q$ where $\mathbf Q_{z,k} = \sum_{i = 1}^n \sum_{l = 1}^n [\mathbf{D}_z]_{i, l} [\mathbf{D}_k]_{i, l}$ for $z, k \in [K]$\\
\textcolor{white}{X} 3. For $t = 1, \dots, T$ iterations (or until convergence): \\
\textcolor{white}{XX} 3.1 For $z=1,\ldots,K$: \\
\textcolor{white}{XXX} 3.1.1 Compute eigendecomposition $\mathbf U \mathbf \Lambda \mathbf U^\top = (\mathbf W_z - \sum_{k \neq z} \mathbf Q_{z,k} \hat \bSigma_k^{(t)})/\mathbf Q_{z,z}$ \\
\textcolor{white}{XXX} 3.1.2 Update $\bSigma_z^{(t+1)} = \mathbf U \mathbf \Lambda^+ \mathbf U^\top$
\end{algorithm}

% \begin{algorithm}\caption{Blockwise coordinate descent algorithm for solving (\ref{eq:mvREHEunpen})} \label{alg:original}
% \setstretch{1.15}
% \textcolor{white}{X} 1. Compute $\mathbf W_G = \sum_{i = 1}^n \sum_{l = 1}^n [\K_n]_{i,l} \Y_{i, :} \Y_{l, :}^\top$ and $\mathbf W_E = \sum_{i = 1}^n  \Y_{i, :} \Y_{i, :}^\top$. \\
% \textcolor{white}{X} 2. Compute $Q_{G, G} = \sum_{i = 1}^n \sum_{l = 1}^n [\K_n]_{i, l}^2$, $Q_{G, E} = \sum_{i = 1}^n [\K_n]_{i, i}^2$, and let $Q_{E, E} = n$.\\
% \textcolor{white}{X} 3. For $t = 1, \dots, T$ iterations or until convergence: \\
% \textcolor{white}{XX} 3.1 Compute eigendecomposition $\mathbf U_G \mathbf \Lambda_G \mathbf U_G^\top = (\mathbf W_G - Q_{G, E} \bSigma_E^{(t)})/Q_{G, G}$ \\
% \textcolor{white}{XX} 3.2
% Update $\bSigma_G^{(t+1)} = \mathbf U_G \mathbf \Lambda_G^+ \mathbf U_G^\top$ \\
% \textcolor{white}{XX} 3.2 Compute eigendecomposition $\mathbf U_E \mathbf \Lambda_E \mathbf U_E^\top = (\mathbf W_E - Q_{G, E} \bSigma_G^{(t)})/Q_{E, E}$ \\
% \textcolor{white}{XX} 3.2
% Update $\bSigma_E^{(t+1)} = \mathbf U_E \mathbf \Lambda_E^+ \mathbf U_E^\top$
% \end{algorithm}

Since the unrestricted mvHE estimator in (\ref{eq:mvHE}) is obtained from a generalized estimating equation for observations $\Y_{i,j}\Y_{l,m}$ ($i,l \in [n], j,m \in [q]$) with an identity working correlation matrix, we can leverage results on generalized estimating equations \parencite{xie2003asymptotics} to establish asymptotic properties of our estimators. Specifically, under mild assumptions on the row covariance matrices $\mathbf{D}_k$, the mvHE estimator is consistent. Additionally, due to this consistency, if the true $\bSigma_k$ are positive definite, the mvHE and mvREHE estimators are asymptotically equivalent and are thus both consistent. We formally state these results in the following proposition, relegating the proof to Appendix \ref{sec:propproof}.
\begin{prop}
Suppose the model in \eqref{eq:genmodel} holds and the row covariance matrices $\mathbf{D}_k$ each satisfy $[\mathbf{D}_k]_{i, l} \leq \rho_h^*, h = |i - l|$, for a sequence $\{\rho_h^*\}_{h = 0}^\infty$ satisfying $\lim_{h \to \infty} \rho_h^* = 0$, and the true column covariance matrices are positive definite: $\bSigma_k \succ 0$. Then the mvHE and mvREHE estimators are consistent. \label{prop:consistency}
\end{prop}
The condition of Proposition~\ref{prop:consistency} is commonly satisfied in genetic studies, where the dataset consists of individuals from many unrelated families, i.e., the genetic relatedness ($\K_n$) and shared household ($\Hn_n$) matrices are approximately block diagonal up to permutation of the rows/columns. Despite the asymptotic equivalence of the mvHE and mvREHE estimators, our simulation results in Section \ref{sec:simulations} show a clear finite sample advantage to using mvREHE. 

% \subsection{Extensions} \label{sec:extensions}

% \subsubsection{High-dimensional traits} \label{sec:highdim}

% In this section, we a consider variant of our method based on dimension reduction that may provide lower estimation error in the setting where $q$ is large relative to $n$. It additionally may provide further computational benefits in the high-dimensional setting. Specifically, let $\Y^r = \Y\mathbf{V}$ be a dimension-reduced version of $\Y$ where $\mathbf V$ is orthogonal. For example, we could set $\mathbf{V}$ to be the first $r$ right singular vectors of $\mathbf{Y}$, so that $\Y^r$ would be the first $r$ classical principal components of $\Y$. Then we model $\Y^r$ as $\Y^r = \sum_{k = 0}^K \boldsymbol{\Gamma}_k^r$ where $\boldsymbol{\Gamma}_k^r \sim \mathcal{MN}(\mathbf{0}_{n \times r}, \mathbf{D}_k, \bSigma_k^r)$. We obtain estimates $\hat \bSigma_k^r$ using our method and define our estimates as $\hat \bSigma_k = \mathbf{V} \hat \bSigma_k^r \mathbf{V}^\top$. 

% This procedure can be motivated by considering the case that the rows of each of the $\boldsymbol{\Gamma}_k$ all lie in a common low-dimensional subspace, so we perform covariance estimation in that subspace rather than in the original space. However, in simulations, we show that this procedure is sometimes beneficial even when this case does not hold.

% Note that while this procedure could be used with other estmators, using mvREHE still has an advantage over other estimators since the value of $r$ can be selected using cross-validation in a computationally efficient manner. Additionally, vanilla mvREHE is simply a special case of this procedure, with $r = q$.

% \subsubsection{Functional traits} \label{sec:functional}

% Often times, neuroimaging traits can be measured along a continuum, though only discrete measurements are obtained. For example, consider the model where the noiseless data is of the form
% $$
% Y_i(t) = \sum_{k = 0}^K \gamma^{(k)}_i(t)
% $$
% where $\gamma^{(k)}_i(t)$ are jointly mean-zero Gaussian random variables with $\cov(\gamma^{(k)}_i(t), \gamma^{(k)}_l(s)) = [\mathbf D^{(k)}]_{i,l} C^{(k)}(t, s)$. In this setting, we are interested in estimating the covariance functions $C^{(k)}$. We obtain noisy measurements of $Y_i$ at $q$ discrete points $t_1, \dots, t_q$, where
% \begin{equation} \label{eq:funcmodel}
% Y_{i,j} = \sum_{k = 0}^K \gamma^{(k)}_{i}(t_j) + \delta_{i,j}
% \end{equation} 
% where $\delta_{i,j} \sim \mathcal{N}(0, \sigma^2)$ are i.i.d. Arranging the random effects in matrix form, where $\boldsymbol{\Gamma}^{(k)}_{i,j} = \gamma^{(k)}_i(t_j)$ and $\boldsymbol \delta_{i,j} = \delta_{i,j}$, 
% we have that $\boldsymbol{\Gamma}^{(k)} \sim \mathcal{MN}(\mathbf{0}_{n \times q}, \mathbf D^{(k)}, \bSigma^{(k)})$ and $\boldsymbol \delta \sim \mathcal{MN}(\mathbf{0}_{n \times q}, \I_n, \sigma^2 \I_q)$, where $[\bSigma^{(k)}]_{i,j} = C^{(k)}(t_i, t_j)$. Without any additional assumptions, with $\mathbf D_0 = \I_n$, it is not possible to identify $\boldsymbol{\Gamma}_0$ and $\boldsymbol{\delta}$, since they both have the same row-covariance, so $\boldsymbol{\Gamma}_0 + \boldsymbol{\delta}$ also follows a matrix-normal distribution with the same row-covariance. However, in the functional data setting described, it is reasonable to assume $C^{(k)}$ are smooth covariance functions, thus it is possible to estimate them even in the presence of noise. Specifically, we propose to fit model (\ref{eq:simplemodel}), but apply a local polynomial smoother to the estimates $\hat \bSigma^{(k)}$ with the diagonal removed, to obtain estimates of $C^{(k)}$.

% \subsection{Relationship with traditional heritability analyses}

% For univariate traits, heritability is often defined as the ratio of the variance of the genetic contribution and the variance of the observed trait. That is, the heritability of the $j$th trait is defined as 
% $$
% h_j^2 = \frac{[\bSigma_G]_{j,j}}{[\bSigma_G]_{j,j} + [\bSigma_C]_{j,j} + [\bSigma_E]_{j,j}}.
% $$

% The notion of \textit{multi-trait heritability proportion} was introduced by \textcite{ge2016multidimensional} by considering the following quantity:
% \begin{equation} \label{eq:heritability}
% \bar h^2 = \frac{\tr(\bSigma_G)}{\tr(\bSigma_G) + \tr(\bSigma_C) + \tr(\bSigma_E)}.
% \end{equation}
% This can be re-expressed as a weighted average of the individual-trait heritabilities, with the weight for each trait being proportional to the overall variance of that trait:
% $$
% \bar{h}^2 = \sum_{j = 1}^q \frac{[\bSigma_G]_{j, j} + [\bSigma_C]_{j, j} + [\bSigma_E]_{j, j}}{\tr(\bSigma_G) + \tr(\bSigma_C) + \tr(\bSigma_E)} h_j^2,
% $$
% % = \sum_{j = 1}^q \frac{[\bSigma_G]_{j, j}}{[\bSigma_G]_{j, j} + [\bSigma_C]_{j, j} + [\bSigma_E]_{j, j}} \frac{[\bSigma_G]_{j, j} + [\bSigma_C]_{j, j} + [\bSigma_E]_{j, j}}{\tr(\bSigma_G) + \tr(\bSigma_C) + \tr(\bSigma_E)} 
% where $h_j^2$ is the individual-trait heritability for the $j$th trait. Therefore, the multi-trait heritability proportion does not capture any of the additional insight that the multivariate model provides with the covariance of traits in each of the components, and it is often estimated using univariate models for each trait separately.

% Both univariate and multivariate heritability only involve the (diagonal entries) of each of the component's covariance matrices, and are thus a byproduct of disentangled multivariate data analysis.

\subsection{Estimation of latent genetic regression parameters} \label{sec:regression_estimation}

Once we obtain an estimate $\hat{\bSigma}_G$, we can define plug-in estimates of the genetic regression parameters. However, rather than directly solving the plug-in version of \eqref{eq:geneticbeta_cov}, the high-dimensional nature of connectome data motivates the use of regularized estimators to improve prediction accuracy and prevent overfitting.

For each functional connection $j$, we consider the general framework where we solve:
\begin{equation} \label{eq:reg_regression}
\hat{\boldsymbol{\beta}}_{G,j} = \argmin_{\boldsymbol{\beta} \in C} \boldsymbol{\beta}^\top \hat{\bSigma}_G^{\mathcal{S},\mathcal{S}} \boldsymbol{\beta} - 2\boldsymbol{\beta}^\top \hat{\bSigma}_G^{\mathcal{S},\mathcal{F}_j} + P(\boldsymbol{\beta})
\end{equation}
where $P(\boldsymbol{\beta})$ is a penalty function and $C$ is a constraint set. This framework accommodates various regularization approaches, and we focus on three specific variants that are well-suited to connectome data.

The first two variants employ standard ridge and lasso penalties: $P(\boldsymbol{\beta}) = \lambda_1 \|\boldsymbol{\beta}\|_2^2$ for ridge and $P(\boldsymbol{\beta}) = \lambda_2 \|\boldsymbol{\beta}\|_1$ for lasso. These approaches address overfitting through shrinkage but do not leverage the matrix structure of brain connectivity. The ridge variant admits a closed-form solution, while the lasso variant can be efficiently solved using standard coordinate descent with active set and warm start optimizations.

The third variant, based on tensor regression, exploits the natural matrix organization of structural connectomes, where connections are arranged by their anatomical origin and destination regions. Rather than treating structural connections as an unstructured vector, this approach leverages their inherent spatial organization.

Specifically, let $\text{mat}(\mathbf{g}_i^{\mathcal{S}}) \in \mathbb{R}^{p \times p}$ denote the matrix representation of the genetic component of the structural connectome, where $p$ is the number of brain regions. For each functional connection $j$, we parameterize the linear predictor as $\langle \boldsymbol{\beta}_{1,j} \boldsymbol{\beta}_{2,j}^\top, \text{mat}(\mathbf{g}_i^{\mathcal{S}}) \rangle_F$, where $\boldsymbol{\beta}_{1,j}, \boldsymbol{\beta}_{2,j} \in \mathbb{R}^{p \times r}$ and $r$ is the rank parameter. This factorization allows connections that share anatomical endpoints to have related regression coefficients, reflecting the spatial organization of brain networks.

To achieve this structure, we define the constraint set as $C = \{\boldsymbol{\beta} : \text{mat}(\boldsymbol{\beta}) = \boldsymbol{\beta}_1 \boldsymbol{\beta}_2^\top, \boldsymbol{\beta}_1, \boldsymbol{\beta}_2 \in \mathbb{R}^{p \times r}\}$, which restricts the matrix form of $\boldsymbol{\beta}$ to have a low-rank factorization. In practice, we optimize over the factors $\boldsymbol{\beta}_{1,j}$ and $\boldsymbol{\beta}_{2,j}$ directly and add penalties on each factor for identifiability and shrinkage:
\begin{equation}\label{eq:tensor_regression}
\hat{\boldsymbol{\beta}}_{1,j}, \hat{\boldsymbol{\beta}}_{2,j} = \argmin_{\boldsymbol{\beta}_1, \boldsymbol{\beta}_2 \in \mathbb{R}^{p \times r}} \vect(\boldsymbol{\beta}_1\boldsymbol{\beta}_2^\top)^\top \hat{\bSigma}_G^{\mathcal{S},\mathcal{S}} \vect(\boldsymbol{\beta}_1\boldsymbol{\beta}_2^\top) - 2\vect(\boldsymbol{\beta}_1\boldsymbol{\beta}_2^\top)^\top \hat{\bSigma}_G^{\mathcal{S},\mathcal{F}_j} + \lambda\left(\|\boldsymbol{\beta}_1\|_F^2 + \|\boldsymbol{\beta}_2\|_F^2\right)
\end{equation}

We solve this optimization problem using alternating minimization: we  update $\boldsymbol{\beta}_{1}$ while holding $\boldsymbol{\beta}_{2}$ fixed, then update $\boldsymbol{\beta}_{2}$ while holding $\boldsymbol{\beta}_{1}$ fixed, and repeat until convergence. Each individual update admits a closed-form solution, making the optimization computationally efficient despite the non-convex nature of the overall problem. We then set $\hat{\boldsymbol{\beta}}_{G,j} = \vect(\hat{\boldsymbol{\beta}}_{1,j}\hat{\boldsymbol{\beta}}_{2,j}^\top)$.


Although we focused on latent genetic regression parameters, these parameters can similarly be defined for the common environmental and unique environmental components.
Moreover, while we focus on linear regression analysis here, the estimate $\hat{\boldsymbol{\Sigma}}_G$ serves as a summary statistic that can be used for many other multivariate data analyses of the genetic component $\mathbf{g}_i$ that are functions of estimated covariance matrices, including, for instance, principal component analysis and canonical correlation analysis.


\vspace{-3.5mm}

% \section{Disentangling genetic and environmental contributions in the human connectome} \label{sec:application}
% \vspace{-3.5mm}

% \subsection{Data}

% In this study, we analyze MRI data from 955 participants in the S1200 HCP dataset, described in Section~\ref{sec:data}. This dataset includes multiple MRI modalities: anatomical, diffusion-weighted, and resting-state functional MRIs. The MRI images were preprocessed using the minimal preprocessing HCP pipeline \parencite{glasser2013minimal}. The resulting anatomical MRI images were used to reconstruct the mid-thickness surface, which represents the geometry of the cortical surface. These surfaces were registered across subjects using the MSM surface matching method \parencite{robinson2014msm}. 

% For each surface, we delineate 68 spatially localized regions of interest (ROIs) using the Desikan-Killiany atlas \parencite{desikan2006automated}, which are consistent across subjects. 

% All 955 participants have kinship information available. Specifically, our dataset includes unrelated subjects, full siblings, half-siblings, dizygotic twins, and monozygotic twins. We encode kinship information in a $955 \times 955$ relatedness matrix, $\K_n$, representing the expected proportion of genome shared between any two subjects due to common ancestry. We compute this matrix from HCP pedigree data using the SOLAR-ECLIPSE imaging genetics software \parencite{koran2014study}, with entries set to $1$ for monozygotic twins, $0.5$ for dizygotic twins or full siblings, $0.25$ for half-siblings, and $0$ for unrelated subjects. The entries of the shared environment matrix $\mathbf{H}_n$ are defined as a binarized version of $\K_n$; specifically $[\mathbf{H}_n]_{i, l} = \mathbbm{1}\left([\mathbf{K}_n]_{i, l} > 0\right)$, based on the assumption that siblings are raised in the same household.



% \vspace{-3.5mm}

% \subsection{Structural connectomes}


% Structural connectomes were constructed from diffusion and structural MRI, as previously described \parencite{kiar2017highthroughput}. In the minimally pre-processed HCP data, a diffusion tensor was estimated in each voxel and streamline estimates of brain connection trajectories were propagated using the Euler delta crossings algorithm, as implemented in DIPY \parencite{garyfallidis2014dipy}. Each subject's anatomy was registered to the MNI152 atlas and for every pair of regions in the \color{red}{Desikan-Killiany atlas}, with merged associated left-right ROIs, a count of the fiber streamlines connecting those two ROIs was produced, resulting in a $68 \times 68$ matrix for each subject, which we refer to as the structural connectome. The subject-level connectomes used here are publicly available in the Neurodata portal \parencite{vogelstein2018communitydeveloped} and through the AWS Open Data program.
% \vspace{-3.5mm}

% \subsection{Functional connectomes}
% The HCP minimally pre-processed fMRI data provide a time-series at each location of the mid-thickness surface describing local neuronal activity. These were acquired in four 15 minute sessions. For each ROI, we computed a spatially averaged time-series describing the average neuronal activity pattern within that ROI using fMRI data from the first session. The resulting time-series were used to compute a $68 \times 68$ correlation matrix between ROIs, which we refer to as the functional connectome.
% \vspace{-3.5mm}


\section{Disentangling genetic and environmental contributions in the human connectome} \label{sec:application}
\vspace{-3.5mm}
\subsection{Data and family structure}

We analyze MRI data from 955 participants in the S1200 HCP dataset, including anatomical, diffusion-weighted, and resting-state functional MRIs preprocessed using the minimal preprocessing HCP pipeline \parencite{glasser2013minimal}. Cortical surfaces are registered across subjects using MSM surface matching \parencite{robinson2014msm}, then parcellated into 68 regions of interest (ROIs) using the Desikan-Killiany atlas \parencite{desikan2006automated}.


Our sample includes participants with varying genetic relatedness: monozygotic twins, dizygotic twins,  full siblings, half-siblings, and unrelated subjects. We encode this family structure using a $955 \times 955$ kinship matrix $\K_n$ representing the expected sharing of the genome, calculated from the HCP pedigree data using SOLAR-ECLIPSE \parencite{koran2014study} with kinship values of 1 for monozygotic twins, 0.5 for dizygotic twins and full siblings, 0.25 for half-siblings, and 0 for unrelated subjects. The shared environment matrix $\mathbf{H}_n$ is defined as $[\mathbf{H}_n]_{i, l} = \mathbbm{1}([\mathbf{K}_n]_{i, l} > 0)$, assuming siblings share the same household.

\vspace{-3.5mm}
\subsection{Connectome construction}

We used structural connectomes that were constructed from diffusion and structural MRI, as previously described \parencite{kiar2017highthroughput}. Using the minimally preprocessed HCP data, diffusion tensors were estimated in each voxel and streamline estimates of brain connection trajectories were propagated using the Euler delta crossings algorithm, as implemented in DiPy \parencite{garyfallidis2014dipy}. Each subject's anatomy was registered to the MNI152 atlas and fiber streamlines connecting each pair of regions in the Desikan-Killiany atlas were counted, resulting in a $68 \times 68$ structural connectome matrix for each subject. The subject-level connectomes used here are publicly available in the Neurodata portal \parencite{vogelstein2018communitydeveloped} and through the AWS Open Data program.



For functional connectomes, the HCP minimally preprocessed fMRI data provide a time-series at each location of the mid-thickness surface describing local neuronal activity. For each ROI, we compute a spatially averaged time-series describing the average neuronal activity pattern within that ROI using fMRI data from the first session. We use the resulting time-series to compute a $68 \times 68$ correlation matrix between ROIs, which we refer to as the functional connectome.


\vspace{-3.5mm}

% \subsection{Statistical analysis}

% \subsubsection{Estimation of covariance matrices} \label{subsec:estimationapplication}

% While numerous representation models for functional and structural connectomes have been proposed in the literature \parencite{arsigny2006logeuclideana, arsigny2006geometric, dryden2009noneuclidean, zhang2019tensor, lila2022functional, zhao2022covarianceoncovariance}, in this work, we adopt a simple approach, commonly used in neuroscience studies for its interpretability \parencite[see, for example,][]{smith2015linking,xia2018linked}. Specifically, we represent each functional and structural connectome by vectorizing the full $68 \times 68$ connectivity matrices, yielding one vector per subject encapsulating the connectomes of the $i$th subject. This vector is 9248-dimensional, as there are 4624 connections from each $68 \times 68$ structural and functional connectivity matrix. We use linear regression to adjust for the effect of age, age squared, sex, intracranial volume, and cortical volume, as in \textcite{colclough2017heritability}, and use the residuals as our mean-zero $\mathbf y_i$.

% While we do not exploit the matrix structure of the connectomes for covariance estimation, we leverage this structure in downstream applications, particularly in the latent genetic tensor regression analysis described in Section~\ref{sec:regression_estimation}.

% To estimate the covariance matrices $\bSigma_k$, we used the proposed mvREHE estimator, since the multivariate REML estimator is not computationally feasible with $q = 9248$ --- we estimate that REML would be intractable at this scale, while the mvREHE estimator remains computationally efficient. To estimate the covariance matrices $\bSigma_k$, we used the proposed mvREHE estimator, since the multivariate REML estimator is not computationally feasible with $q = 4556$ --- we estimate that REML would be intractable at this scale, while the mvREHE estimator remains computationally efficient. Before applying the proposed mvREHE estimator, we first rescale the columns of $\Y$ such that they have unit standard deviation. We then rescale the final estimates to be on the original scale. This is equivalent to reweighting the loss function \eqref{eq:mvREHEunpen} to down-weight the contribution of connections that have a larger scale and up-weight the contribution of connections that have a smaller scale; this reweighting ensures that the estimation focuses on the correlation structure rather than being dominated by connections with larger raw variances. The reweighted estimator can also be shown to be consistent for the true parameter.


\subsection{Statistical analysis}

\subsubsection{Estimation of covariance matrices} \label{subsec:estimationapplication}

While numerous representation models for functional and structural connectomes have been proposed in the literature \parencite{arsigny2006logeuclideana, arsigny2006geometric, dryden2009noneuclidean, zhang2019tensor, lila2022functional, zhao2022covarianceoncovariance}, we adopt a simple vectorization approach commonly used in neuroscience studies for its interpretability \parencite[see, for example,][]{smith2015linking,xia2018linked}. We represent each functional and structural connectome by vectorizing the full $68 \times 68$ connectivity matrices, yielding a 9,248-dimensional vector per subject (4,624 connections from each structural and functional matrix). Although this representation introduces redundancy due to symmetry in the connectivity matrices, it facilitates the tensor regression analysis without incurring a computational penalty in the covariance estimation step, as the covariance estimation runtime is dominated by the sample size rather than the feature dimension (see Proposition \ref{prop:runtime}). We use linear regression to adjust for the effects of age, age squared, sex, intracranial volume, and cortical volume, as in \textcite{colclough2017heritability}, and use the residuals as our mean-zero $\mathbf y_i$. While we do not enforce the matrix structure of connectomes at the stage of covariance estimation, we leverage this structure in the downstream latent genetic tensor regression analysis described in Section~\ref{sec:regression_estimation}.


We estimate the covariance matrices $\bSigma_k$ using our proposed mvREHE estimator, since multivariate REML is computationally intractable with $q = 9,248$ dimensions. Before applying mvREHE, we rescale the columns of $\Y$ to have unit standard deviation, then rescale the final estimates back to the original scale. This is equivalent to reweighting the loss function \eqref{eq:mvREHEunpen}, ensuring that estimation focuses on the correlation structure rather than being dominated by connections with larger raw variances. The reweighted estimator remains consistent for the true parameter.


\vspace{-3.5mm}

\subsubsection{Latent genetic and environmental regression analysis of structure-function coupling}

\textcite{gu2021heritability} previously investigated the genetics of structure-function coupling by first estimating the correlation between structural and functional connections for each subject---finding a moderate association with values up to 0.4--and then performing univariate heritability analysis on these resulting coefficients. That is, they estimate how much of the observed within-subject structure-function coupling can be attributed to genetic factors.

This moderate within-subject coupling does not translate to the population level, where the structural connectome shows virtually no linear predictive power for functional connections across subjects (Section \ref{subsec:datanalaysisresults}). In contrast, we hypothesize that a fundamental, genetically-encoded association between structure and function exists at the population level, but that it is obscured by unique environmental exposures.

Therefore, rather than asking, ``\textit{how heritable is an individual's summary of structure-function coupling?}'', as in \textcite{gu2021heritability}, we ask, ``\textit{what is the population-level linear relationship between structural and functional connectomes that is encoded in genetics?}'' By separating the genetic component from the environmental component, our methodology aims to uncover the structure-function coupling that can be inherited. This population-level framework also allows us to model more complex connectivity patterns in which a functional connection is influenced by multiple structural connections, which heritability analyses of subject-specific structure–function summaries cannot capture.

To study standardized effects, we apply the latent genetic regression framework described in Section~\ref{sec:regression_estimation}, but estimate the genetic regression parameters in \eqref{eq:geneticbeta} defined with $\tilde{\mathbf{g}}_i = \diag(\bSigma_G)^{-1/2} \mathbf{g}_i$ instead of $\mathbf{g}_i$. In practice, this means using the estimated genetic correlation matrix rather than the estimated genetic covariance matrix in the regularized loss function. This yields regression coefficients that represent the number of standard deviations the genetic component of functional connectivity changes for each standard deviation change in the genetic component of structural connectivity, holding the genetic component of all other structural connections fixed.

In our application to HCP data, we consider all three regularization methods (ridge, lasso, and tensor regression). To select the optimal regularization parameter and tuning parameters, we estimate the genetic correlation matrix and regression parameters on one half of the data and compute the genetic-$R^2$ on the other half, averaging the estimates of out-of-sample genetic-$R^2$ across splits, then choose the regularization method and tuning parameters that maximize this averaged out-of-sample genetic $R^2$. This approach allows us to determine which regularization strategy best captures the genetic structure-function relationships for each connection while avoiding overfitting.

We perform the same regression analysis on the common environment and unique environment components to compare linear predictability across latent components. This enables us to determine whether structure-function couplings are primarily driven by genetic factors, environmental factors, or a combination of both. We also perform classical linear regression using the observed data to compare against traditional analyses of connectome data and to assess how much the genetic decomposition reveals about structure-function coupling that would otherwise be masked in the observed data.

We report the out-of-sample genetic $R^2$ values for the selected method and tuning parameters. The resulting genetic $R^2$ values quantify the proportion of variance in each functional connection's genetic component that can be explained by a linear function of the genetic components of all structural connections. High genetic $R^2$ values indicate strong genetically-encoded structure-function coupling, while low values suggest little evidence of a linear relationship between the genetic components of structural and functional connectivity.


% Again, $\boldsymbol{\beta}_G^*$ can be rewritten in terms of the covariance matrices:
% $$
% \boldsymbol{\beta}_G^* = \left[\left\{\diag([\bSigma_G])^{-1/2}\bSigma_G\diag([\bSigma_G])^{-1/2}\right\}_{\mathcal{S}, \mathcal{S}}\right]^{-1} \left\{\diag([\bSigma_G])^{-1/2}\bSigma_G\diag([\bSigma_G])^{-1/2}\right\}_{\mathcal{S}, \mathcal{F}}
% $$ 

% We are also interested in quantifying the linear predictability of each of the functional connections based on the structural connections. The proportion of variance explained in the $j$th functional connection by the structural connections for a specific coefficient vector $\boldsymbol \beta$ is given by
% $$
% R^2_{G, j}(\boldsymbol{\beta}) = 1- \frac{SSE_{G, j}(\boldsymbol{\beta})}{SST_{G, j}(\boldsymbol{\beta})} = 1 - \frac{\E \left[ \| [ \mathbf{\tilde g}_{i}]_{\mathcal{F}_j} - \boldsymbol \beta_{:,j}^\top [ \mathbf{\tilde g}_{i}]_{\mathcal{S}} \|_2^2 \right]}{\E \left[ \| [ \mathbf{\tilde g}_{i}]_{\mathcal{F}_j} - \E[[\mathbf{\tilde g}]_{\mathcal{F}_j}] \|_2^2 \right]} = 1 - \frac{[\bSigma_G]_{\mathcal{F}_j, \mathcal{F}_j} - 2 \boldsymbol \beta_{:,j}^\top [\bSigma_G]_{S, \mathcal{F}_j} + \boldsymbol \beta_{:,j}^\top [\bSigma_G]_{\mathcal{S}, \mathcal{S}} \boldsymbol \beta_{:,j} }{[\bSigma_G]_{\mathcal{F}_j, \mathcal{F}_j}}.
% $$
% Once again, this parameter only depends on the genetic covariance matrix.


% :
% $$
% \boldsymbol{\hat \beta}_{G, j} = \left[\left\{\diag([\hat \bSigma_G])^{-1/2} \hat \bSigma_G\diag([\hat \bSigma_G])^{-1/2}\right\}_{\mathcal{S}, \mathcal{S}} + \lambda_j \I_{|\mathcal{S}|} \right]^{-1} \left\{\diag([\hat \bSigma_G])^{-1/2} \hat \bSigma_G\diag([\hat \bSigma_G])^{-1/2}\right\}_{\mathcal{S}, \mathcal{F}_j},
% $$
% where $\lambda_j$ is a tuning parameter we select by minimizing the cross-validated $\hat{SSE}_{G, j}({\boldsymbol{\hat \beta}})$. Note that the cross-validation procedure requires estimating the covariance twice in each of the cross-validation folds (once for the training set and once for the test set), making this estimator impractical with existing inefficient covariance estimation algorithms.


% \subsubsection{Genetic modes of variation analysis}

% Prior studies have reported that genetics contributes strongly to connectomes \parencite[][, which we also replicate this in Appendix \ref{subsec:h2}]{adhikari2018comparison, barber2021heritability, ge2017heritability}. We aim to further characterize the shape of these genetic contributions. We perform an exploratory analysis similar to principal components analysis, to identify the direction that explains the most variation in the disentangled genetic component of the connectome. Due to differences in scale between the different connections and between the structural and functional connectomes, we first standardize by the genetic standard deviation. Specifically, we are interested in estimating and interpreting the first eigenvector of $\diag(\bSigma_G)^{-1/2} \bSigma_G \diag(\bSigma_G)^{-1/2}$; that is, the main mode of variability as captured by the genetic correlation structure. The other modes of variation are given by the subsequent eigenvectors. To estimate these directions, we apply our methodology to get an estimate of the genetic covariance $\hat \bSigma_G$, and use this plug-in estimate to determine the directions that explain the most variation in the genetic component.

% %$$
% %\argmin_{\mathbf w} \E \left[ \|\diag(\bSigma_G)^{-1/2} %\G_i - \mathbf{w} \mathbf{w}^\top %\diag(\bSigma_G)^{-1/2} \G_i\|_2^2 \right],
% %$$ which can be shown to be 


% If we were to simply perform classical principal components analysis on the observed data, we would obtain an estimate of the direction that explains the most variation in the sum of genetic, common environment, and unique environment components instead. That is, we would be estimating the first eigenvector of $\diag(\bSigma_O)^{-1/2} \bSigma_O \diag(\bSigma_O)^{-1/2}$, where $\bSigma_O = \bSigma_G + \bSigma_C + \bSigma_E$. However, this approach does not isolate the genetic contribution from environmental ones in general, and thus it does not characterize the shape of the genetic contributions individually.

\subsection{Results} \label{subsec:datanalaysisresults}

\subsubsection{Genetic and environmental regression analysis of the structure-function coupling}


Figure~\ref{fig:r2} shows the out-of-sample estimates of the proportion of variance of each functional connection that can be explained by all structural connections for the observed data (observed $R^2$) and each latent component (genetic $R^2$, common environment $R^2$, and unique environment $R^2$). Each $R^2$ value represents the best performance across all three regularized latent regression methods and their respective tuning parameters for that functional connection. The results reveal striking differences in structure-function predictability across observed and latent components.


\begin{figure}
\centering
\includegraphics[width=\textwidth]{Figures/Revision/r2.pdf}
\caption{Proportion of variance explained ($R^2$) in functional connectivity by structural connectivity for each functional connection. Results are shown for (a) observed data and decomposed latent components: (b) genetic, (c) common environment, and (d) unique environment. Each cell represents the out-of-sample $R^2$ for predicting one functional connection using all structural connections as predictors. Darker shades of red indicate higher predictability.}\label{fig:r2}
\end{figure}


In the observed data, structural connections show minimal linear predictive power for functional connections, with most $R^2$ values near zero and a maximum of only 0.03. However, when restricted to latent components, a clear pattern emerges. The genetic component shows the strongest structure-function coupling, with many functional connections exhibiting substantial predictability from structural connections (maximum $R^2$ = 0.34). The common environment component shows moderate predictability (maximum $R^2$ = 0.17), while the unique environment component shows virtually no linear relationship between structural and functional connectivity (maximum $R^2$ = 0.04).

These results suggest that the apparent lack of structure-function coupling in observed data masks a strong genetically-encoded relationship that is obscured by unique environmental factors affecting each individual differently.

Different mechanisms could explain the observed phenomenon of high predictability within the genetic component \parencite{vanrheenen2019genetic}, such as horizontal pleiotropy, where a shared set of genes contributes directly or through intermediate endophenotypes to both structural and functional connectivity. Alternatively, a causal relationship could exist, i.e., vertical pleiotropy, although the observed high predictability may also be due to spurious pleiotropy, such as biases in the imaging pipelines used to produce surrogate estimates of functional and structural connectivity.

Figure~\ref{fig:coefs} shows the estimated genetic regression coefficients from all three regularization methods for a functional connection within the default mode network (DFM), which has the highest genetic $R^2$ (0.34). The three methods achieved different out-of-sample genetic $R^2$ values: ridge regression ($R^2$ = 0.34), lasso regression ($R^2$ = 0.27), and tensor regression ($R^2$ = 0.26). We display results from all three methods to contrast their regularization effects.

Since the structural connectome is symmetric, we handle the coefficient visualization differently across methods for interpretability. For ridge and lasso (panels a-b), we sum coefficients from symmetric pairs of structural connections since only their sum affects prediction, displaying results in the lower triangle with upper triangle set to zero. For tensor regression (panel c), we preserve the estimated coefficient matrix to illustrate the low-rank structure.

The coefficient patterns reflect each method's regularization properties: ridge produces distributed shrinkage across connections, lasso yields sparse selection, and tensor regression captures structured spatial relationships through its rank constraint. Interestingly, ridge regression achieves the highest performance for this DFM connection despite tensor regression's incorporation of the structural connectome's matrix structure, suggesting that distributed shrinkage may be more beneficial than structural constraints for this particular connection. 

Across all functional connections, tensor regression is optimal for 44.8\% of connections, ridge regression for 44.2\%, and lasso regression for 11.0\%, indicating that no single regularization approach dominates and that the optimal method depends on the specific functional connection being modeled. These differences demonstrate how regularization choice affects both the performance and interpretability of estimated structure-function coupling in the genetic component.






% $$
% R^2 = 1 - \frac{\E \left[ \| \diag([\bSigma_G]_{\mathcal{Y}, \mathcal{Y}})^{-1/2} \G_{i\mathcal{Y}} - \boldsymbol \beta^\top \diag([\bSigma_G]_{\mathcal{X}, \mathcal{X}})^{-1/2} \G_{i\mathcal{X}} \|_2^2 \right]}{\E \left[ \| \diag([\bSigma_G]_{\mathcal{Y}, \mathcal{Y}})^{-1/2} \G_{i\mathcal{Y}}\|_2^2 \right]}
% $$


\begin{figure}
\centering
\includegraphics[width=\textwidth]{Figures/Revision/beta_G.pdf}
\caption{Estimated genetic regression coefficients for predicting a functional connection within the default mode network (DFM) using all structural connections as predictors. Results shown for (a) tensor regression, (b) lasso regression, and (c) ridge regression. For ridge and lasso, coefficients from symmetric structural connection pairs are summed and displayed in the lower triangle. For tensor regression, the full coefficient matrix structure is preserved to illustrate the low-rank structure.}\label{fig:coefs}
\end{figure}



% \subsubsection{Genetic modes of variation findings}


% % $$
% % \argmin_{\mathbf w} \E \left[ \|\diag(\bSigma_O)^{-1/2} \Y_i - \mathbf w \mathbf w^\top \diag(\bSigma_O)^{-1/2} \Y_i\|_2^2 \right],
% % $$

% %   In the structural connectome, increases in the structural connection between the salience and visual and salience and medial-parietal regions are associated with decreases in the majority of other structural connections, while they are associated with increases in the majority of functional connections. However, increases in the structural connection between the salience and visual and salience and medial-parietal regions are associated with decreases in the rest of the functional connections from the salience region to other regions. 

% We compare the first two genetic and observed primary directions of variation in Figure~\ref{fig:pcs}. This first direction of variability of the observed data captures a uniform change in all functional connections and a corresponding variation in structural connectivity in the same direction. The second mode of variation highlights more pronounced variability in the functional connections within the subnetwork defined by DAT, CCP, and FPT, with structural connections varying in the opposite direction.


% The primary direction of variation driven by the genetic contribution also shows a uniform increase/decrease in all functional connections, suggesting that this mode of inter-subject variability could be partially rooted in genetics. This is accompanied by a variation in structural connectivity indicating that these changes in functional connectivity are associated with opposite changes in the structural connections involving the salience network, particularly between the SAL network and AUD cortex connections.

% The second mode of genetic variation appears to capture a contrasting behavior between the strength of functional connections involving the SAL network, especially those between the SAL network and the SMH/SMM cortices, and all other functional connections. Structural connections also exhibit contrasting behavior between SAL-VIS connections and most other connections. The SAL network is implicated in both estimated modes of genetic variation. Interestingly, our previous analyses on individual-trait heritability proportions (Figure~\ref{fig:h2}a) found that connections involving the SAL region were among the most heritable. Although the functions of the SAL network remain uncertain \parencite{seeley2019salience}, this suggests that its function and interaction with other networks might be non-trivially modulated by genetic factors.


% \begin{figure}[!t]
% \centering
% \includegraphics[width=\textwidth]{Figures/pc_loadings_new.pdf}
% \caption{(a-b) First two dominant modes of variation of the observed structural and functional connectomes. (c-d) The first two genetically driven modes of variation of the structural and functional connectomes. Both genetic modes of variation appear to implicate connections involving the SAL network.}\label{fig:pcs}
% \end{figure}

% Finally, we look at the cumulative proportion of variation explained in each of the components by an increasing number of modes of variation in Figure~\ref{fig:pcsprop} of the Supplementary Materials. 
% % $$
% % 1 - \frac{\E \left[ \|\diag(\bSigma_G)^{-1/2} \G_i - \mathbf{W} \mathbf{W}^\top \diag(\bSigma_G)^{-1/2} \G_i\|_2^2 \right]}{\E \left[ \|\diag(\bSigma_G)^{-1/2} \G_i\|_2^2 \right]}
% % $$
% We see that fewer modes of variation are required to describe the genetic and common environment components as compared to the unique environment component or observed data, indicating that genetics and common environment result in contributions that are more structured than the unique environment contributions. Additionally, this suggests that the observed connectome may look unstructured due to unique environmental exposures. Therefore, analyzing the genetic component alone through its estimated covariance matrix can help provide clues into the underlying structure of the connectome prior to any environmental exposures.






% While numerous representation models for connectomes have been proposed in the literature, in this work we adopt the simplest approach, which is commonly used in neuroscience studies \parencite[see, for example,][]{smith2015linking}. Specifically, we represent each functional and structural connectome using vectors derived from the vectorized lower triangular elements, thereby yielding a matrix where the $i$th row is composed of the representation coordinates of the connectomes of the $i$th subject. We normalize the sub-matrices corresponding to the functional and structural connectomes to each have Frobenius norm 1.

% We compute an SVD of representations and utilize the first $r$ left singular vectors from the functional and structural connectomes to apply the proposed multidimensional heritability model. The principal modes of genetic and environmental variability are computed through eigenanalysis of the genetic and environmental covariance matrices, and are then back-projected onto the original feature space using the right singular vectors of the functional and structural connectomes.

% \subsection{Sensitivity analysis}

% To select $r$, the number of left singular vectors of the functional and structural connectomes, we do a sensitivity analysis to determine the value of $r$ which results in the most stable estimates of genetic and environmental principal directions. Specifically, for each value of $r \in \{10, 25, 50, 75, 100\}$, we estimate the genetic and environmental principal direction from the entire data. We then repeatedly sample 80\% of the data and compute the cosine similarity between the genetic and environmental principal direction estimated from the subsampled data and that computed from the entire data. The results of this sensitivity analysis are shown in Figure~\ref{fig:reproducibility}. Based on these results, we choose to use $r = 50$, since it overall maximizes the average cosine similarity.

% \begin{figure}[!htb]
% \centering
% \includegraphics[width=0.8\textwidth]{Figures/stability.pdf}
% \caption{Stability of the estimated modes of genetic and environmental variation}\label{fig:reproducibility}
% \end{figure}

% \subsection{Results}
% The first and second principal genetic and environmental directions of variation are shown in Figures~\ref{fig:pc1}-\ref{fig:pc2}. We also display the principal directions of variation using classical PCA, referred to as the ``raw PCs''. For the first principal direction, we see that using classical PCA, the connection between almost all pairs of regions in the brain change in the same way along the principal direction of variation. However, the genetic and environmental primary directions of variation are much more complex and biologically meaningful. For example, increased functional connectivity between the cingulo-opercular and auditory regions correlates with reduced connectivity within the auditory region itself, and vice-versa, in the first genetic direction of variation.



% \begin{figure}
% \centering
% \includegraphics[width=\textwidth]{Figures/newpc_2loadings.pdf}
% \caption{Mean and direction of second most variability in community-level connectivity obtained using raw, genetic, and environmental covariance estimates.}\label{fig:pc2}
% \end{figure}

\vspace{-3.5mm}

\section{Simulations} \label{sec:simulations}

In this section, we consider a simulation study based on a related analysis of the HCP data, in an effort to validate our choice of methodology. That is, we simulate data from the additive genetics and common environment model in \eqref{eq:multivariatevariateobservedmodel}, where the row-covariance and column-covariance matrices are set based on estimates from a coarser-resolution analysis of the HCP data.

We group the 68 regions into 10 functional communities and analyze the connectivity at this resolution. For structural connectomes, we sum all structural connections within each community to obtain inter-community connection strengths. For functional connectomes, we average the time series of all regions within each community and compute correlations on these averaged time series. We then apply the mvREHE estimator to estimate the genetic, common environment, and unique environment covariance matrices, which serve as the true parameters for our simulation model.

The genetic relatedness matrix $\K_n$ is defined as follows: if the sample size $n \leq 1,000$, $\K_n$ is the top left $n \times n$ block of the HCP kinship matrix; if $n > 1,000$, then $\K_n$ is an $n \times n$ block diagonal matrix with blocks defined by the $1,000 \times 1,000$ HCP kinship matrix. As in Section \ref{subsec:estimationapplication}, we set the common household matrix $\Hn_n$ to be a binarized version of $\K_n$.

In each replicate, we sample observed data from the model in \eqref{eq:multivariatevariateobservedmodel} and estimate $\bSigma_k$ using mvHE and mvREHE. We do not use multivariate REML as it is computationally intractable with $q = 200$ dimensions. For comparison with existing univariate approaches, we also estimate the diagonal elements of $\bSigma_k$ using univariate HE, REHE, and REML estimators applied separately to each trait. We perform 50 replications for each simulation setting.






% We record runtime, estimation error in $\bSigma_k$ (as measured by spectral error $\|\hat \bSigma_k - \bSigma_k\|_2$), and estimation error in the individual-trait heritability proportions (as measured by L2-error on the vector of heritability proportion of all traits $\|\hat h^2 - h^2 \|_2$). We repeat this 50 times and report the mean runtime and errors.

In Figure~\ref{fig:simdatacombined}a, we compare the runtime of multivariate and univariate estimators. The multivariate methods mvHE and mvREHE show virtually identical computational times, with execution times under a couple of seconds even at $n = 8,000$. The positive definite constraint in mvREHE adds no computational burden due to our efficient algorithm. Both multivariate methods are orders of magnitude faster than multiple univariate estimators, particularly REML, while providing full covariance matrices rather than just diagonal elements.

We next compare the estimation error of the entire covariance matrix using mvHE and mvREHE, the only computationally feasible multivariate estimators in this setting. Figure~\ref{fig:simdatacombined}b shows results for the genetic component, with results for other components in Figure~\ref{fig:simdataspectral} in the Supplementary Materials. The positive semi-definite constraint in mvREHE yields considerably lower estimation error than the truncation-based mvHE, particularly at smaller sample sizes, while having nearly identical runtime performance.


Finally, we compare the prediction performance of latent genetic regression analysis using mvHE and mvREHE for the plug-in covariance estimates. We evaluate all combinations of covariance estimation methods (mvHE vs mvREHE) and regularization approaches (ridge, lasso, tensor regression). Figure~\ref{fig:simdatacombined}c shows the average ratio of achieved genetic-$R^2$ to the oracle genetic-$R^2$ across all functional connections, where the oracle represents the theoretical maximum achievable with the true covariance matrix. Values closer to 1 indicate better recovery of the true predictive relationships in the genetic component. The results demonstrate that mvREHE consistently outperforms mvHE across all regularization methods, with the performance gap particularly pronounced at  larger sample sizes. Results for the common environmental and unique environmental components are shown in Figure~\ref{fig:simdatabeta}. This confirms that the improved covariance estimation accuracy of mvREHE translates directly into better downstream prediction performance across all latent components.

\begin{figure}[!t]
\centering
\includegraphics[width=\textwidth]{Figures/Revision/simulation_data.pdf}
\caption{Performance of the proposed estimation framework across different analyses. (a) Runtime comparison of multivariate methods (mvHE, mvREHE) versus univariate methods (HE, REHE, REML), (b) spectral-norm error of estimated genetic covariance matrix $\hat{\bSigma}_G$ comparing mvHE and mvREHE, and (c) prediction performance of latent genetic regression analysis showing the average ratio of achieved genetic-$R^2$ to oracle genetic-$R^2$ across functional connections for different combinations of covariance estimation methods and regularization approaches, all with varying sample size $n$ and $q = 200$ fixed.}\label{fig:simdatacombined}
\end{figure}

%In Figure~\ref{fig:remlzero} in the Appendix, we compare the estimated and true heritability proportions from a representative replicate of the simulation study. Recall that in our data analysis, univariate REML estimated the heritability of several traits to be zero, and in this simulation we observe it does so for traits which truly have nonzero heritability proportions. On the other hand, the mvREHE estimates are generally close to the true heritability proportions. These simulation results, directly mimicking the observed data, suggest that heritability estimates from mvREHE are more trustworthy than univariate REML.

% Finally, we show the Euclidean-norm of the error of the disentangled regression coefficients for the association between VIS-COP functional connection and the structural connectome, in Figure~\ref{fig:simdatacombined}e for the genetic component and Figure~\ref{fig:simdatabeta} of the Supplementary Materials for the other components. We see that the plug-in estimate of the regression coefficients on the genetic and unique environmental mechanisms have the smallest error using mvREHE, while for the common environment mechanism mvHE and mvREHE perform comparably.

As it was not computationally feasible to directly compare mvREHE with the multivariate REML estimator with $q = 200$, we perform additional low-dimensional simulations $(q \leq 10)$ in Section \ref{subsec:lowdim} of the Supplementary Materials. In Figure \ref{fig:sim1combined}, we observe that the estimation performance of mvREHE nearly exactly matches that of multivariate REML, both for estimation of the entire covariance matrix as well as heritability proportions, while being nearly three orders of magnitude faster than the when $q = 5$ and five orders of magnitude faster when $q = 10$. %Additionally, by extrapolating the runtime complexity of multivariate REML, in Figure \ref{fig:runtime}, we estimate that mvREHE is ten orders of magnitude faster than the multivariate REML estimator in the setting of the data application ($q = 182$ and $n \approx 1000$), reducing runtime from an estimated 300 years to just seconds.

In summary, these simulation results suggest that the mvREHE estimator provides the best of both worlds in terms of computational efficiency and estimation accuracy in the data-inspired settings considered.

% \subsection{High-dimensional traits}

% In this section, we consider a setting with $q = 1000$ outcome variables where the true $\bSigma_k$ have different rates of eigenvalue decay. Specifically, we construct $\bSigma_k = \mathbf{U} \mathbf{D} \mathbf{U}^\top$ where $\mathbf{U} \in \mathbf{R}^{q \times q}$ is an orthogonal matrix and $\mathbf{D}$ is a diagonal matrix with entires $\mathbf{D}_{j, j} = (1/j)^\alpha$ where $\alpha = 1.25, 1, 0.75$ for the ``fast,'' ``moderate,'' and ``slow'' eigenvalue decay settings, respectively. 

% We compare mvHE and mvREHE to dimension-reduced versions of mvREHE and multivariate REML as described in \ref{sec:highdim}. For the dimension-reduction approaches, we set $\mathbf{V}$ to be the first $r$ right singular vectors of $\mathbf{Y}$, where the degree of dimension reduction $r$ is chosen by cross validation for mvREHE and fixed at $r = 5$ for multivariate REML due to computational feasibility.

% \begin{figure}
% \centering
% \includegraphics[width=\textwidth]{Figures/sim2/simulation_figure_spectral_error_n.pdf}
% \caption{Estimation error in the setting of ``fast,'' ``moderate,'' and ``slow'' eigenvalue decay of the true $\bSigma_k$, with $q = 1000$ and varying $n$.}\label{fig:sim2error}
% \end{figure}

% In Figure~\ref{fig:sim2error}, we show the spectral error of estimating $\bSigma_k$ for each of these methods. Vanilla mvREHE greatly improves over mvHE across all settings, and the dimension-reduction variant of mvREHE slightly improves on vanilla mvREHE with small sample sizes and slower eigenvalue decay. As the sample size increases, the error of the mvHE, mvREHE, and cross-validated dimension-reduced mvREHE estimators gets smaller, demonstrating consistency of these methods, but the error of the dimension-reduction version of mvREML levels off, due to bias even with large sample sizes caused by the extreme dimension reduction. Thus, in practice, the improved computational efficiency of mvREHE helps to improve statistical performance as well.

% In Figure~\ref{fig:sim2angle} in Appendix \ref{sec:suppfigs}, we show the angle between the estimated and true main modes of variability in each component. We see that the dimension-reduced version of mvREHE performs the best in all settings, giving slight improvements over vanilla mvREHE when the sample size is small, and large improvements over both mvHE and the dimension-reduced version of multivariate REML across all settings.


% \subsection {Functional traits}

% In this section, we simulate data from the model in (\ref{eq:funcmodel}) with covariance functions defined by
% $$
% C_G(s, t) = C_E(s, t) = \sum_{k = 1}^{50} k^{-2\alpha}\cos(k \pi s)\cos(k \pi t).
% $$
% We consider $\alpha = 1$ and $\alpha = 2$, where the covariance function with $\alpha = 2$ is smoother than that with $\alpha = 1$. We set the variance of the observational noise $\sigma^2 = 1$. We observe the function at $q = 100$ timepoints equally spaced between $0$ and $1$. 

% To estimate the covariance functions, we define the mvHE and mvREHE estimates as piecewise-constant functions given by the elements of the estimated $\hat \bSigma_k$. We compare these to smoothed versions of these estimators, as described in Section \ref{sec:functional}. To get smooth estimates $\hat C_k$, we apply a local linear smoother to the estimated $\hat \bSigma_k$ after removing the diagonal elements, with a Epanechnikov kernel and bandwidth chosen by generalized cross-validation. 

% In Figure~\ref{fig:sim3error}, we show the integrated squared error of estimating $C_k$. We see that with low sample sizes, smoothing helps in estimation of the covariance functions for all components, and mvREHE improves over mvHE for both the smoothed and unsmoothed versions. For estimation of $C_E$, smoothing helps even with large sample sizes, since it is able to recover the smooth covariance function which does not contain the observational noise along the diagonal term, as described in Section \ref{sec:functional}. The unsmoothed estimates of $C_E$ are not consistent due to the lack of identifiability without smoothness.

% \begin{figure}
% \centering
% \includegraphics[width=\textwidth]{Figures/sim3/simulation_figure_squared_error_n.pdf}
% \caption{Estimation error with two covariance functions with different degrees of smoothness, where larger $\alpha$ means smoother. The number of discrete observations along the function is $q = 100$ with $n$ varying.}\label{fig:sim3error}
% \end{figure}

% % Suppose for a moment that  each of the $\bSigma_k$ are  low-rank (meaning the $\Gamma_k$ are low-rank), such that $Y$ is also low-rank. Then we can find an orthogonal matrix $V \in \mathbb{R}^{q \times r}$ such that $Y = YVV^\top$ and $\Gamma_k = \Gamma_k V V^\top$. Let $\Gamma_k^r = \Gamma_k V$. Then $\Gamma_{ki} = V \Gamma_{ki}^r$ so $\bSigma_k = V \bSigma_k^r V^\top$. 

% % While the above setting may not hold in practice, we use this as motivation to reduce the number of parameters to estimate. We consider performance of estimating $\bSigma_k$ under three settings of eigenvalue decay of $\bSigma_k$. 


% % In this setting, [WHY?] we may do linear dimension reduction using a matrix $V \in \mathbb{R}^{q \times r}$ to get $Y^r = YV$ such that $\Gamma_k V V^\top \approx \Gamma_k V V^\top$ where $Y^r V^\top \approx Y$. Then $V^\top \Gamma_{ki}^r \approx \Gamma_{ki}$, so $ \approx $


% % In this case, we propose to estimate the column covariances of the components of $Y^r$

\vspace{-3.5mm}

% \section{Discussion}\label{sec:discussion}

% In this paper, we propose a novel multidimensional analysis of structural and functional connectomes that can dissect the effects of genetic and environmental factors. Our regression analysis characterizing the relationship between functional and structural connectomes restricted to genetic mechanisms extends previous work \parencite{gu2021heritability}, which models structure-function coupling using a single subject-specific coefficient and assumes that each structural pair influences only its corresponding functional pair. By contrast, our analysis allows for more complex global interaction patterns between structural and functional connectomes, revealing a high degree of coherence in structure-function coupling at a genetic level.

% Since the desired regression analysis is not computationally feasible using existing methods, we introduce a computationally efficient method for variance component estimation in multivariate models. This method is based on a restricted generalized method-of-moments estimator, which can be computed with operations scaling quadratically with the product of the sample size and number of traits followed by each iteration scaling cubically with the number of traits.
% The proposed method extends well beyond neuroscience applications, enabling the analysis of complex, high-dimensional data that were infeasible to analyze using existing variance component estimation methods. We demonstrate that the covariance estimates from our multivariate method enable estimates of heritability proportions that are more accurate than the common practice of using REML on each trait individually. Furthermore, by estimating the entire covariance matrix, our method enables a diverse set of multivariate data analyses on the latent components, such as principal components analysis and linear regression used in our analysis of the human structural and functional connectomes. The central role of covariance matrices in multivariate data analysis ensures the broad applicability of our method to various settings, where the observed data is the sum of several latent independent components, yet the primary interest is in analyzing the components separately.


\section{Discussion}\label{sec:discussion}

We developed a novel latent genetic regression analysis framework that characterizes relationships between genetic components of brain connectivity. Unlike previous approaches that estimate subject-specific structure-function coupling and then assess its heritability \parencite{gu2021heritability}, our method directly models genetic contributions, additionally allowing for a population-level understanding of complex interaction patterns where functional connections can be influenced by multiple structural connections.

The computational feasibility of this analysis required methodological innovation in multivariate variance component estimation. Existing estimators like multivariate REML are computationally intractable for high-dimensional connectome data. We introduce a constrained method-of-moments estimator that guarantees positive semi-definite covariance estimates. Simulation studies demonstrate that our block coordinate descent algorithm efficiently incorporates the positive semi-definite constraint while maintaining practically identical runtime as the unconstrained estimator, and that the method nearly matches multivariate REML's estimation accuracy while providing dramatic computational speedups.

Our analysis reveals that functional connectivity is substantially more predictable from structural connectivity when restricted to genetic components (maximum $R^2$ = 0.34) compared to observed data, which show minimal predictive relationships (maximum $R^2$ = 0.03). This finding suggests that unique environmental factors mask genetically-encoded relationships between structural and functional connectomes.

The proposed method extends well beyond neuroscience applications, applying to any multivariate setting where observed data represents the sum of latent genetic and environmental components. This enables the analysis of complex, high-dimensional phenotypes which can be modelled as the sum of independent latent components, which was previously computationally infeasible. By estimating entire covariance matrices rather than just variances, our method enables a diverse set of multivariate analyses restricted to individual latent components, including principal component analysis, canonical correlation analysis, and the regression framework demonstrated here. These analyses fundamentally rely on covariance structure.


\section{Code availability}

An open-source R package implementing the proposed estimator is available at 
% \url{https://github.com/[REDACTED FOR DOUBLE BLINDNESS]}.
\url{https://github.com/keshav-motwani/mvREHE}.

% \setcounter{maxnames}{10}
%\newpage 

\printbibliography


\newpage 
\spacingset{1}


\appendix
\appendixpage

\section{Proof of Lemma 1} \label{prooflemma}

\begin{proof}
Expressing the Frobenius norm as the trace of a quadratic form, using linearity of the trace operator, and applying an eigendecomposition gives us that
\begin{align*}
\hat \bSigma &= \argmin_{\bSigma \succeq 0}  \sum_{w = 1}^W \|\mathbf{S}_w - a_w \bSigma \|_F^2 \\
&= \argmin_{\bSigma \succeq 0}   \left\|\left(\sum_{w = 1}^W a_w^2\right)^{-1} \sum_{w = 1}^W a_w \mathbf{S}_w - \bSigma \right\|_F^2\\
&=\argmin_{\bSigma \succeq 0}   \left\|\mathbf S - \bSigma \right\|_F^2 \\
&=\argmin_{\bSigma \succeq 0}   \left\|\mathbf U \mathbf \Lambda \mathbf U^\top  - \bSigma \right\|_F^2 \\
&=\argmin_{\bSigma \succeq 0}   \left\|\mathbf \Lambda  - \mathbf U ^\top \bSigma \mathbf U \right\|_F^2.
\end{align*}
Since $\mathbf \Lambda$ is diagonal, we can restrict ourselves to the case where $\mathbf U ^\top \bSigma \mathbf U$ is also diagonal. Note that the optimal $\bSigma$ is PSD if and only if $\mathbf U ^\top \bSigma \mathbf U$ is PSD, since $\mathbf U$ is an orthogonal matrix. Thus, the diagonal elements of $\mathbf U ^\top \bSigma \mathbf U$ must be non-negative. If $\lambda_j \geq 0$, then the best case is that the $j$th diagonal element of $\mathbf U ^\top \bSigma \mathbf U$ is also $\lambda_j$. If $\lambda_j < 0$, then the best case is that the $j$th diagonal element of $\mathbf U ^\top \bSigma \mathbf U$ is $0$. Hence, the minimizer is achieved by $\hat \bSigma = \mathbf U \mathbf \Lambda^+ \mathbf U^\top$. 
\end{proof}

\newpage

\section{Proof of Proposition \ref{prop:consistency}} \label{sec:propproof}
\begin{proof}
First, we define some notation. Let $\dot Y = \vect(\Y^\top)$. Then under the model in equation (\ref{eq:genmodel}), $\dot Y \sim \mathcal{N}(0, \sum_{k = 0}^K \mathbf D_k \otimes \bSigma_k)$. Let $\mathbf E_{jl}$ is a $q$ by $q$ matrix with 1 in the $j,l$ and $l,j$ entries and 0 elsewhere. Then equation (\ref{eq:mvHE}) can be written as the solution to
\begin{align}\label{eq:rewrite}
&\argmin_{\bSigma_k} \| \dot \Y \dot \Y^\top - \sum_{k = 0}^K \mathbf D_k \otimes \bSigma_k \|_F^2 \\
&= \argmin_{\bSigma_k} \| \dot \Y \dot \Y^\top -  \sum_{k = 0}^K \sum_{j \geq l} [\bSigma_k]_{j, l} (\mathbf D_k \otimes \mathbf E_{jl}) \|_F^2 \\
&= \argmin_{\sigma} \| \tilde \Y - \tilde{\mathbf{X}} \boldsymbol \sigma \|_2^2
\end{align}
where $\boldsymbol \sigma_{g(k, j, l)} := [\bSigma_k]_{j,l}$, $\tilde \Y = \vect(\dot \Y \dot \Y^\top)$, and $\tilde{\mathbf{X}}_{:, g(k, j, l)} = \vect(\mathbf D_k \otimes \mathbf E_{jl})$ and $g$ is a function which linearizes the indices $k \in \{0\} \cup [K]$, $j, l \in [q]$ with $j \geq l$: $$g(k, j, l) = \frac{kq(q+1)}{2} + (l - 1)(q + 1) - \frac{l(l - 1)}{2} + j - (l - 1).$$ Therefore, (\ref{eq:rewrite}) is the ordinary least squares estimate of $\sigma$ in the linear model
$$
\tilde \Y = \tilde{\mathbf{X}} \boldsymbol \sigma + \boldsymbol \delta
$$
where $\cov(\boldsymbol \delta) = \mathbf R$. We can see that the pre-truncation mvHE estimator is the special case of a generalized estimating equation estimator with the identity matrix as the working correlation matrix. Thus, as in \textcite{yue2021rehe}, we use results from \textcite{xie2003asymptotics} on asymptotics of generalized estimating equations. 

One can verify that Assumption 1 implies that the correlation of $\boldsymbol \delta$ is bounded by $|\mathbf R_{a, b} / \sqrt{\mathbf R_{a, a} \mathbf R_{b, b}}| < \rho_h$ with $h = |a - b|$ for a sequence $\rho_h$ with $\lim_{h \to \infty} \rho_h = 0$, since 
\begin{align*}
\cov(\Y_{i, j}\Y_{l, m}, \Y_{i', j'}\Y_{l', m'}) &= E(\Y_{i, j}\Y_{l, m}\Y_{i', j'}\Y_{l', m'}) - E(\Y_{i, j}\Y_{l, m})E(\Y_{i', j'}\Y_{l', m'}) \\ 
&= E(\Y_{i, j}\Y_{i', j'})E(\Y_{l, m}\Y_{l', m'}) + E(\Y_{i, j}\Y_{l', m'}) E(\Y_{l, m}\Y_{i', j'}) \\
&= \cov(\Y_{i, j},\Y_{i', j'})\cov(\Y_{l, m},\Y_{l', m'}) + \cov(\Y_{i, j},\Y_{l', m'}) \cov(\Y_{l, m},\Y_{i', j'})
\end{align*}
and 
$$
\cov(\Y_{i, j}, \Y_{l, m}) = \mathbb{E}(\Y_{i, j} \Y_{l, m}) = [\bSigma_E]_{j, m} [\I_n]_{i, l} + [\bSigma_G]_{j, m} [\K_n]_{i, l}.
$$
Therefore, under this assumption, the pre-truncation mvHE estimator is consistent \parencite{xie2003asymptotics, yue2021rehe}. Additionally, the truncated mvHE estimator is consistent by the continuous mapping theorem.

% \begin{theorem}
% (Consistency) Under Assumption 1, as $n \to \infty$, $\hat \bSigma_{\text{mvHE/mvREHE}} \overset{p}{\to} \bSigma$
% \end{theorem}

% If $\K_n$ is block diagonal with blocks $\K_n^{m}$, then we additionally get asymptotic normality of the mvHE estimator. 

% \begin{assumption}
% (Strong Condition, Assumption 2 of \textcite{yue2021rehe}) The covariance matrix $\K_n$ is block diagonal with blocks $\K_n^{(1)}, \dots, \K_n^{(M)}$.
% \end{assumption}

% Let $Y = (Y^{(1)\top}, \dots, Y^{(M)\top})^\top$ where $Y^{(m)}$ is defined to be the subvector of $Y$ corresponding to the $m$th block of $\K_n$. Let $\dot Y^{(m)} = \vect(Y^{(m)})$ and $\tilde Y^{(m)} = \vect(\dot Y^{(m)} \dot Y^{(m)\top})$. Letting $\tilde Y_S = (\tilde Y^{(1)\top}, \dots, \tilde Y^{(M)\top})^\top$ and 
% $\tilde X_{S{:, g(k, j, l)}} = (\vect(D_k^{(1)} \otimes E_{jl})^\top, \dots, \vect(D_k^{(M)} \otimes E_{jl})^\top)^\top$, under Assumption 2, we have that 
% \begin{align*}
% & \argmin_{\bSigma} \| \tilde Y - \tilde X \bSigma \|_2^2
% = \argmin_{\bSigma} \| \tilde Y_S - \tilde X_S \bSigma \|_2^2
% \end{align*}
% Additionally, since $\I_n = I_n$ and $Y$ is matrix normal, this implies independence of the $Y^{(m)}$'s and thus of the $\tilde Y^{(m)}$'s. 

% Letting $R^{(m)} = \cor(Y^{(m)})$ and
% $$
% W_M = \sum_{m = 1}^M \tilde X^{(m)\top} R^{(m)} X^{(m)},
% $$
% $$
% H_M = \sum_{m = 1}^M \tilde X^{(m)\top} X^{(m)},
% $$
% we have that
% $$
% W_M^{-1/2} H_M (\hat \bSigma_{mvHE} - \bSigma) \overset{d}{\to} \mathcal{N}(0, I)
% $$
% as long as the maximum block size does not go to infinity too fast \parencite{xie2003asymptotics, yue2021rehe}. 

Now we argue asymptotic equivalence of the mvHE and mvREHE estimators if the true $\bSigma_k$ are positive definite. Note that the mvREHE estimate is the same as the mvHE estimate if the covariance matrices estimated by mvHE are PSD. By openness of the space of positive definite matrices, there exists an open ball of positive definite matrices around the true $\bSigma_k$ which the mvHE estimate will asymptotically belong to, by consistency. Therefore, the mvREHE estimate will asymptotically be the same as the the mvHE estimate, and will also be consistent. 
\end{proof}

% \section{Motivation for dimension reduction extension}

% Suppose the features in the observations in each of the components are restricted to some low-dimensional subspace of $\mathbb{R}^q$, described by the column space of some orthogonal matrix $\mathbf V$. That is, $\boldsymbol{\Gamma}_k = \boldsymbol{\Gamma}_k^r \mathbf{V}^\top$ for all $k$. Then the observations in $\Y$ are also restricted to this subspace, that is, $\Y = \Y^r \mathbf V^\top$, so $\Y^r = \Y \mathbf{V}$. Thus
% $$
% \Y^r = \Y \mathbf V = \sum_{k = 0}^K \boldsymbol{\Gamma}_k \mathbf V = \sum_{k = 0}^K \boldsymbol{\Gamma}_k^r
% $$
% where $\boldsymbol{\Gamma}_k^r \sim \mathcal{MN}(\mathbf{0}_{n \times q}, \mathbf{D}_k, \bSigma_k^r)$
% and $\bSigma_k^r \equiv \mathbf V ^\top \bSigma_k \mathbf{V}$. Thus we can get an estimate $\bSigma_k^r$ using $\Y \mathbf V$ as our data and recover $\bSigma_k$ as $ \bSigma_k = \mathbf{V} \bSigma_k^r \mathbf V ^\top$. In practice, $\mathbf V$ is unknown, so we estimate this with SVD as well and use a plug-in estimate.

% \begin{align*}
% &\argmin_{\bSigma_k \succeq 0}  \sum_{i = 1}^n \sum_{l = 1}^{n} \left\|\Y_{i,:}^\top \Y_{l,:}   - \sum_{k = 0}^K [\mathbf{D}_k]_{i, l} \mathbf V \bSigma_k^r \mathbf V^\top \right\|_F^2 \\
% &\argmin_{\bSigma_k \succeq 0}  \sum_{i = 1}^n \sum_{l = 1}^{n} \left\|\Y_{i,:}^\top \Y_{l,:}   - \mathbf V\sum_{k = 0}^K [\mathbf{D}_k]_{i, l}  \bSigma_k^r \mathbf V^\top \right\|_F^2 \\
% &= \argmin_{\bSigma_k \succeq 0}  \sum_{i = 1}^n \sum_{l = 1}^{n} -2\tr\left(\Y_{i,:}^\top \Y_{l,:} \mathbf V \sum_{k = 0}^K [\mathbf{D}_k]_{i, l}  \bSigma_k^r \mathbf V^\top \right) + \tr\left(\mathbf V \sum_{k = 0}^K [\mathbf{D}_k]_{i, l}  \bSigma_k^r \mathbf V^\top \mathbf V \sum_{k = 0}^K [\mathbf{D}_k]_{i, l}  \bSigma_k^r \mathbf V^\top \right)\\
% &= \argmin_{\bSigma_k \succeq 0}  \sum_{i = 1}^n \sum_{l = 1}^{n} -2\tr\left(\mathbf V^\top \Y_{i,:}^\top \Y_{l,:} \mathbf V  \sum_{k = 0}^K [\mathbf{D}_k]_{i, l} \bSigma_k^r  \right) + \tr\left\{\left(\sum_{k = 0}^K [\mathbf{D}_k]_{i, l}  \bSigma_k^r  \right)^\top \left(\sum_{k = 0}^K [\mathbf{D}_k]_{i, l}  \bSigma_k^r  \right)\right\}\\
% \end{align*}

\renewcommand{\thefigure}{S\arabic{figure}}



% \section{Supplemental Data Analysis Results}



% \subsection{Heritability of the human connectome} \label{subsec:h2}

% We now use our estimates of the covariance matrices to estimate the individual-trait heritability proportions for both structural and functional connections between regions in the human brain, to see whether there is an advantage to using a multivariate estimator even when not strictly required. We focus on individual-trait heritability rather than multi-trait heritability because, as mentioned previously, the latter is simply a weighted average of the former, and there are large differences in scale between traits in our study. In Figure~\ref{fig:h2}a, we display the heritability proportion estimates obtained using our method for each of the structural and functional connections. We see that the heritability estimates of some structural and functional connections are as much as 23\%, with the most heritable structural connections involving the salience and frontal/medial parietal regions and the most heritable functional connections involving the visual, medial parietal, and ventral attention regions. 

% Since individual-trait heritability proportions can also be estimated using univariate models, we compare the estimates with those obtained using univariate REML on each of the traits separately in Figure~\ref{fig:h2}b. We observe that for many connections in the brain, REML estimates a zero or near-zero heritability proportion, while mvREHE estimates a much larger heritability proportion. REML also estimates a high heritability proportion of nearly 60\% in one case. In the Supplementary Material, we show that REML is susceptible to extreme estimates and that the mvREHE estimates of individual-trait heritability are more accurate in this setting, showing the advantage of a multivariate modelling approach. Additionally, the mvREHE estimator for the entire covariance matrix is more than 20 times faster than the REML estimator for just the diagonal elements.

% \begin{figure}
% \centering
% \includegraphics[width=\textwidth]{Figures/heritability.pdf}
% \caption{(a) Individual-trait heritability of each structural and functional connection. (b) Comparison of heritability estimates between mvREHE and univariate REML. We observe that REML often estimates
% a zero or near-zero heritability proportion for many brain connections, and in one case, as high as 60\%. Such extreme values are likely artifacts, as demonstrated by simulations. Moreover, computing the mvREHE estimate for the entire covariance matrix is over 20 times faster than calculating REML estimates for just the diagonal elements.}\label{fig:h2}
% \end{figure}


\section{Supplemental Simulation Results} 
\label{sec:suppfigs}

\subsection{Data-inspired simulation (described in Section \ref{sec:simulations} of main text)}


\begin{figure}[!htb]
\centering
\includegraphics[width=\textwidth]{Figures/Revision/spectral.pdf}
\caption{Estimation error of all three covariance matrices $\bSigma_k$ (genetic, common environmental, and unique environmental components) comparing mvHE and mvREHE estimators with varying sample size $n$. Error is measured using the spectral norm of the difference between estimated and true covariance matrices. Results demonstrate the superior performance of mvREHE over mvHE across all components, particularly at smaller sample sizes.}
\label{fig:simdataspectral}
\end{figure}



\begin{figure}[!htb]
\centering
\includegraphics[width=\textwidth]{Figures/Revision/r2_prop.pdf}
\caption{Prediction performance of latent regression analysis for common environmental and unique environmental components. Results show the average ratio of achieved $R^2$ to oracle $R^2$ across all functional connections for different combinations of covariance estimation methods (mvHE vs mvREHE) and regularization approaches (ridge, lasso, tensor regression) with varying sample size $n$.}
\label{fig:simdatabeta}
\end{figure}

% \begin{figure}[h]
% \centering
% \includegraphics[width=0.8\textwidth]{Figures/simdata/simulated_h2.pdf}
% \caption{Estimated individual-trait heritability proportions for one replicate of data-inspired simulation. Univariate REML is susceptible to to erroneously estimating heritability proportions to be zero for traits which are truly heritable.}\label{fig:mvremlfail}
% \end{figure}


\subsection{Low-dimensional simulation} \label{subsec:lowdim}

In this section, we perform simulations in the low-dimensional setting, where it is possible to directly compare with the multivariate REML estimator. Due to its statistical efficiency, the multivariate REML estimator is commonly considered a gold standard, despite being computationally inefficient. We set the covariance matrices $\bSigma_k$ to be scaled unstructured correlation matrices with $q = 5, 10, 20$ traits. The scaling is defined such that the multi-trait heritability matches that of the estimated model from the real data application. The unstructured correlation matrices are fixed to be a uniform draw from the space of correlation matrices, across all simulation replicates. We consider sample sizes ranging from 250 to 8000.

% Finally, we scale the true column covariance matrices such that the proportion of variance explained by the genetic component (heritability proportion) is 10\% and that explained by the common environment component is 10\%, as defined in (\ref{eq:heritability}), in order to roughly match the data application.

\begin{figure}
\centering
\includegraphics[width=\textwidth]{Figures/sim1/simulation_figure_combined_n.pdf}
\caption{(a) Runtime, (b) spectral-norm error of $\hat \bSigma_G$, and (c) Euclidian-norm error of the individual-trait heritability proportions $h^2$ in the setting of unstructured, low-dimensional $\bSigma_k$ with varying $q$ and $n$.}\label{fig:sim1combined}
\end{figure}

In Figure~\ref{fig:sim1combined}a, we show the runtime for each of the estimators considered. First, we can see that mvREHE has a similar runtime to mvHE --- that is, imposing the positive semi-definite constraint does not affect the runtime in these settings. We can also see that mvREHE is nearly three orders of magnitude faster than the multivariate version of REML when $q = 5$ and nearly five orders of magnitude faster when $q = 10$. With large sample sizes and $q = 10$ or $q = 20$ with any sample size, the multivariate REML estimator cannot be computed in reasonable time; thus, these results are not shown. We can see that increasing the number of traits $q$ does not have a large impact on the runtime of mvREHE, while for multivariate REML the increase in runtime is substantial. Additionally, the multivariate REML estimator is susceptible to convergence issues, resulting in a large number of replicates returning no estimates for this estimator (over 90\% of replicates failing in some settings with large sample sizes; see Figure~\ref{fig:mvremlfail} in the Appendix). Finally, the mvREHE estimator is even faster than fitting separate univariate models for each trait using any of the univariate estimators, while estimating the entire covariance matrix rather than the diagonal only.


Given that the mvREHE estimator is the most computationally efficient, one may ask if it compromises on estimation performance. In Figure~\ref{fig:sim1combined}b, we show the spectral error of estimating the covariance matrix of the genetic component, with results for the other components given in Figure~\ref{fig:sim1spectral}. We can see that the PSD constraint in the mvREHE estimator results in a large improvement over the mvHE estimator, nearly exactly matching the estimation accuracy of the multivariate REML estimator in all cases where the REML estimate could be computed. 

We also compare estimation performance of the individual-trait heritability proportions in Figure~\ref{fig:sim1combined}c, as measured by the Euclidian-norm error on the vector of heritability proportion of all traits. Since the heritability proportion depends only on the diagonal elements of the covariance matrices, we compare against the univariate estimators as well. We again see that mvREHE and multivariate REML have nearly identical performance, and that these estimators greatly outperform their commonly-used univariate counterparts. Finally, we see a large benefit from using a PSD constraint in mvREHE rather than simply truncating the estimate in mvHE. 

% In order to better understand the runtime complexity of the multivariate REML estimator, as well as estimate the expected runtime of the multivariate REML estimator in the setting of the data application ($q = 182$ and $n \approx 1000$), we repeat the same simulations keeping $n = 1000$ fixed and varying $q$ between 3 and 10 for multivariate REML, and between 3 and 182 for the other methods. We fit a function of the form $C q^\alpha$ to the relationship between $q$ and the runtime, where $\alpha$ and $C$ are estimated from the observed runtimes for different values of $q$ across 50 replicates. The results are shown in Figure \ref{fig:runtime}. We observe that the observed total runtime of the multivariate REML estimator with the sample size $n$ fixed roughly scales with $q^5$, whereas mvREHE, mvHE, and the univariate methods seem to scale linearly with $q$. Additionally, we estimate the runtime of multivariate REML to be over 300 years in the setting of the data application, with $q = 182$ and $n = 1000$, compared to just seconds for the mvREHE estimator. 

In Figure~\ref{fig:sim1combined2component}, we present the results of a simulation akin to the one described in this section, but with only the genetic and unique environment components, in order to facilitate comparison with the optimized REML algorithm proposed by \textcite{zhou2014efficient} which only applies to two component models. In the two component setting, all methods have very similar estimation performance. While the algorithm proposed by \textcite{zhou2014efficient} is up to two orders of magnitude faster than the \texttt{lme4}-based implementation of REML with $q = 10$, it is still two orders of magnitude slower than mvREHE when $q = 10$ and four orders of magnitude slower than mvREHE when $q = 20$. 



\setcounter{figure}{0}

\begin{figure}[!htb]
\centering
\includegraphics[width=\textwidth]{Figures/sim1/simulation_figure_spectral_error_n.pdf}
\caption{Estimation error of all three covariance matrices $\bSigma_k$ in the setting of low-dimensional $\bSigma_k$
with varying sample size $n$.}
\label{fig:sim1spectral}
\end{figure}

\begin{figure}[h]
\centering
\includegraphics[width=0.8\textwidth]{Figures/sim1/simulation_figure_na_prop_n.pdf}
\caption{Proportion of simulation replicates in low-dimensional simulation setting in which the multivariate REML estimator fails to converge. Note that no results are reported for the settings with $q=10$ and $n = 4000$ and $8000$, as REML estimation is computationally prohibitive in these cases.}\label{fig:mvremlfail}
\end{figure}

% \begin{figure}[h]
% \centering
% \includegraphics[width=0.8\textwidth]{Figures/sim1/simulation_figure_time_q.pdf}
% \caption{Runtime in the setting of unstructured $\bSigma_k$ with $n = 1000$ fixed and $q$ varying. Fitted lines approximate runtime as $C q^\alpha$, where $\alpha$ and $C$ are estimated from the observed runtime over 50 replicates. On the log-log scale shown, $\alpha$ is the slope of the line. These results suggest that fitting the multivariate REML estimator in the setting of the data application would take approximately $10^{10}$ seconds, or over 300 years.}\label{fig:runtime}
% \end{figure}

\begin{figure}[h]
\centering
\includegraphics[width=\textwidth]{Figures/sim1-2comp/simulation_figure_combined_n.pdf}
\caption{(a) Runtime, (b) spectral-norm error of $\hat \bSigma_G$, and (c) Euclidian-norm error of the individual-trait heritability proportions $h^2$ in the setting of unstructured, low-dimensional $\bSigma_k$ with varying $q$ and $n$ with only the genetic and unique environmental components (two component model).  GEMMA is the optimized algorithm for computing the multivariate REML estimator, proposed by \textcite{zhou2014efficient}. }\label{fig:sim1combined2component}
\end{figure}


% \begin{figure}
% \centering
% \includegraphics[width=0.8\textwidth]{Figures/sim2/simulation_figure_max_principal_angle_1_n.pdf}
% \caption{Angle between estimated and true main modes of variation in the setting of ``fast,'' ``moderate,'' and ``slow'' eigenvalue decay of the true $\bSigma_k$, with $q = 1000$ and varying $n$.}\label{fig:sim2angle}
% \end{figure}


\section{Extension to function-valued traits} \label{sec:functional}

Often times, neuroimaging traits can be measured along a continuum, though only discrete measurements are obtained. Here, we show how the proposed mvREHE estimator can be easily combined with ideas from functional data analysis \parencite{yao2005functional} to incorporate smoothness into the covariance estimation step, similar to \textcite{risk2021ace}.

Consider the following model for the noiseless data of the $i$th subject:
$$
y_i(t) = g_i(t) + c_i(t) + e_i(t)
$$
for $t \in [0, 1]$, where $g_i$, $c_i$, and $e_i$ are mean-zero Gaussian processes with smooth covariance functions $C_G$, $C_C$, and $C_E$ respectively. Suppose we observe noisy versions of $y_i$ at $q$ discrete points $t_1, \dots, t_q$, denoted as the vector $\mathbf{y}_i$, where
\begin{equation} \label{eq:funcmodel}
[\mathbf{y}_i]_j = y_i(t_j) + \delta_{ij} = g_i(t_j) + c_i(t_j) + e_i(t_j) + \delta_{ij}
\end{equation} 
for $j = 1, \dots, q$, where $\delta_{ij} \sim \mathcal{N}(0, \sigma^2)$ are i.i.d. Then,
\begin{equation}  \label{eq:multivariateindivmodel}
\mathbf{y}_i \overset{d}{=} \mathbf{g}_i + \mathbf{c}_i + \mathbf{e}_i + \boldsymbol{\delta}_i
\end{equation}
where
$\mathbf{g}_i \sim \mathcal{N}(0, \bSigma_G)$, $\mathbf{c}_i \sim \mathcal{N}(0, \bSigma_C)$, $\mathbf{e}_i \sim \mathcal{N}(0, \bSigma_E)$, and $\boldsymbol{\delta}_i \sim \mathcal{N}(0, \sigma^2 \I_q)$, and $[\bSigma_k]_{j,l} = C_k(t_j, t_l)$ for $k \in \{G, C, E\}$. 

Let $\mathbf{Y}, \mathbf{G}, \mathbf{C}, \mathbf{E}$, and $\boldsymbol{\Delta}$ be matrices with rows $\mathbf{y}_i$,  $\mathbf{g}_i$, $\mathbf{c}_i$, $\mathbf{e}_i$, and $\boldsymbol{\delta}_i$, respectively. As in the observed data model in equation \eqref{eq:multivariatevariateobservedmodel}, we assume that $\G \sim \mathcal{MN}(\mathbf{0}_{n \times q}, \K_n, \bSigma_G)$, $\C \sim \mathcal{MN}(\mathbf{0}_{n \times q}, \Hn_n, \bSigma_C)$, $\E \sim \mathcal{MN}(\mathbf{0}_{n \times q}, \I_n, \bSigma_E)$, and $\boldsymbol{\Delta} \sim \mathcal{MN}(\mathbf{0}_{n \times q}, \I_n, \sigma^2 \I_q)$. Without any additional assumptions, with $\mathbf D_0 = \I_n$, it is not possible to identify $\E$ and $\boldsymbol{\Delta}$, since they both have the same row-covariance, therefore $\E + \boldsymbol{\Delta}$ also follows a matrix-normal distribution with the same row-covariance. However, in the functional data setting described, it is reasonable to assume $C_{k}$ are smooth covariance functions, thus it is possible to estimate them even in the presence of noise. Specifically, we propose to fit model \eqref{eq:multivariatevariateobservedmodel}, but, similar to \parencite{yao2005functional, risk2021ace}, apply a 2D local polynomial smoother to the estimates $\hat \bSigma_{k}$ with the diagonal removed, to obtain estimates of $C_{k}$.

\subsection{Simulation results}

In this section, we simulate data from the model in equation (\ref{eq:funcmodel}) with covariance functions defined by
$$
C_G(s, t) = C_E(s, t) = \sum_{k = 1}^{50} k^{-2\alpha}\cos(k \pi s)\cos(k \pi t).
$$
We consider $\alpha = 1$ and $\alpha = 2$, where the covariance function with $\alpha = 2$ is smoother than that with $\alpha = 1$. We set the variance of the observational noise $\sigma^2 = 1$. We sample the function at $q = 100$ timepoints equally spaced between $0$ and $1$. 

To estimate the covariance functions, we define the mvHE and mvREHE estimates as piecewise-constant functions given by the elements of the estimated $\hat \bSigma_k$. We compare these to smoothed versions of these estimators, as described in Section \ref{sec:functional}. To get smooth estimates $\hat C_k$, we apply a local linear smoother to the estimated $\hat \bSigma_k$ after removing the diagonal elements, with an Epanechnikov kernel and bandwidth chosen by generalized cross-validation. 

In Figure~\ref{fig:sim3error}, we show the integrated squared error of estimating $C_k$. We see that with low sample sizes, smoothing helps in estimation of the covariance functions for all components, and mvREHE improves over mvHE for both the smoothed and unsmoothed versions. For estimation of $C_E$, smoothing helps even with large sample sizes, since it is able to recover the smooth covariance function which does not contain the observational noise along the diagonal term, as described in Section \ref{sec:functional}. Note that the unsmoothed estimates of $C_E$ are not consistent due to the lack of identifiability without smoothness. This highlights the versatility of the proposed framework and its ability to be combined with other covariance estimation approaches to incorporate additional information about the data.

\begin{figure}
\centering
\includegraphics[width=\textwidth]{Figures/sim3/simulation_figure_squared_error_n.pdf}
\caption{Estimation error with two covariance functions with different degrees of smoothness, where larger $\alpha$ means smoother. The number of discrete observations along the function is $q = 100$ with $n$ varying.}\label{fig:sim3error}
\end{figure}

\end{document}